\documentclass[11pt,a4paper]{report}
% Compilation : tectonic Cours_SER.tex
% Figures à données : python3 tools/make_figures.py (appelle Code/ser et Code/aero)
% Préambule commun (repris du cours sur le modèle aéro)
\usepackage{iftex}
\ifPDFTeX
  \usepackage[utf8]{inputenc}
  \usepackage[T1]{fontenc}
  \usepackage{lmodern}
\else
  \usepackage{fontspec}
\fi
\usepackage[french]{babel}
\usepackage[a4paper,margin=2.3cm]{geometry}
\usepackage{amsmath,amssymb,mathtools}
\usepackage{xcolor}
\usepackage{tikz}
\usetikzlibrary{arrows.meta,calc,positioning,shapes.geometric,decorations.pathreplacing,decorations.markings,patterns,babel}
\usepackage{pgfplots}
\pgfplotsset{compat=1.18}
\usepgfplotslibrary{groupplots}
\usepackage{tabularx,booktabs,array}
\usepackage{enumitem}
\setlist{nosep,leftmargin=*}
\usepackage{listings}
\usepackage[font=small,labelfont=bf]{caption}
\usepackage[raggedright]{titlesec}
\usepackage[most]{tcolorbox}
\usepackage[colorlinks=true,linkcolor=blue!50!black,urlcolor=blue!50!black,citecolor=blue!50!black]{hyperref}

\definecolor{cblue}{HTML}{2A78D6}
\definecolor{corange}{HTML}{EB6834}
\definecolor{cgreen}{HTML}{1BAF7A}
\definecolor{cink}{HTML}{1F2328}
\definecolor{cmuted}{HTML}{6B6F76}

% ---------------------------------------------------------------- notations
\newcommand{\dd}{\mathrm{d}}
\newcommand{\Ue}{U_e}
\renewcommand{\Re}{\mathit{Re}}
\newcommand{\vect}[1]{\mathbf{#1}}
\newcommand{\code}[1]{\texttt{#1}}

% ---------------------------------------------------------------- encadrés
\tcbset{coursbox/.style={enhanced,breakable,boxrule=0pt,leftrule=2.5pt,arc=0pt,
  left=6pt,right=6pt,top=4pt,bottom=4pt,fonttitle=\bfseries\small,coltitle=cink,
  attach title to upper={\ }}}
\newtcolorbox{definition}[1][]{coursbox,colback=cblue!5,colframe=cblue,title={Définition#1.}}
\newtcolorbox{resultat}[1][]{coursbox,colback=cgreen!6,colframe=cgreen,title={À retenir#1.}}
\newtcolorbox{projet}[1][]{coursbox,colback=corange!6,colframe=corange,title={Dans notre projet#1.}}
\newtcolorbox{piege}[1][]{coursbox,colback=cmuted!7,colframe=cmuted,title={Piège#1.}}
\newtcolorbox{codebox}[1][]{coursbox,colback=cmuted!3,colframe=cblue!60,title={Dans le code#1.}}

% ---------------------------------------------------------------- exercices
\newcounter{exo}[chapter]
\renewcommand{\theexo}{\thechapter.\arabic{exo}}
\newenvironment{exercice}[1][]{\refstepcounter{exo}\par\medskip\noindent
  \textbf{\color{cblue}Exercice~\theexo}\if\relax\detokenize{#1}\relax\else\ (#1)\fi.\ }{\par}
\newcommand{\exercices}{\section*{Exercices}\addcontentsline{toc}{section}{Exercices}
  \noindent\textit{Corrigés en annexe~\ref{annexe:corriges}. Chercher d'abord sans les regarder.}}

\tikzset{
  flow/.style={draw=cink,rounded corners=2pt,fill=white,align=center,font=\small,inner sep=4pt},
  arr/.style={-{Stealth[length=2.2mm]},cink,thick},
}

\lstset{basicstyle=\ttfamily\footnotesize,breaklines=true,columns=fullflexible,frame=single,rulecolor=\color{cmuted!50},
  backgroundcolor=\color{cmuted!6},keywordstyle=\color{cblue}\bfseries,commentstyle=\color{cmuted}\itshape,
  stringstyle=\color{corange},showstringspaces=false,language=Python,
  literate={é}{\'e}1{è}{\`e}1{à}{\`a}1{ê}{\^e}1{ç}{\c c}1{ô}{\^o}1{î}{\^i}1{ù}{\`u}1{É}{\'E}1{→}{$\to$}1}

\setcounter{secnumdepth}{2}
\setcounter{tocdepth}{1}
\setlength{\parindent}{0pt}
\setlength{\emergencystretch}{3em}
\setlength{\parskip}{0.45em}


\definecolor{cviolet}{HTML}{7A4FB3}
\newtcolorbox{pistes}[1][]{coursbox,colback=cviolet!6,colframe=cviolet,title={Pistes d'amélioration#1.}}
\newcommand{\dB}{\,\mathrm{dB}}

\newcommand{\shat}{\hat{\vect s}}

% ============================================================================
% Itérations du projet HackathonQuantique : couleurs et marqueurs communs.
% Inclus par tous les cours et par la fiche projet, pour qu'une itération ait
% la même couleur partout. Pour ajouter une itération : une couleur, un nom,
% une date, une ligne dans \legendeiterations.
% ============================================================================
\definecolor{itV0}{HTML}{8D6E63}   % brun   : V0, modèle de la note d'intention
\definecolor{itV1}{HTML}{C2185B}   % framboise : V1, chaîne complète conforme à la note
\definecolor{itV2}{HTML}{00838F}   % cyan foncé : V2, problème reformulé
\definecolor{itV3}{HTML}{E09A00}   % ambre  : V3, concurrent équitable et problème plus grand
\definecolor{itV4}{HTML}{5E35B1}   % indigo : V4, réservé

\makeatletter
\newcommand{\itername}[1]{\@nameuse{itname@#1}}
\newcommand{\iterdate}[1]{\@nameuse{itdate@#1}}
\newcommand{\defiter}[3]{\@namedef{itname@#1}{#2}\@namedef{itdate@#1}{#3}}
\makeatother
\defiter{V0}{Modèle de la note d'intention}{24--30 oct. 2026}
\defiter{V1}{Chaîne complète conforme à la note}{4--5 nov. 2026}
\defiter{V2}{Problème reformulé}{5--6 nov. 2026}
\defiter{V3}{Concurrent équitable, problème plus grand}{à partir du 6 nov. 2026}

% Pastille : \iterbadge{V1}
\newcommand{\iterbadge}[1]{\tikz[baseline=(b.base)]\node[fill=it#1,text=white,rounded corners=2pt,
  inner xsep=3pt,inner ysep=1.5pt,font=\sffamily\bfseries\scriptsize] (b) {#1};}

% Bandeau de début de chapitre : \iterbanner{V0}
\newcommand{\iterbanner}[1]{\par\noindent\begin{tikzpicture}
  \node[fill=it#1!12,draw=it#1,line width=0.6pt,rounded corners=2pt,text width=\dimexpr\linewidth-12pt\relax,
    inner sep=4pt,anchor=west,font=\small] at (0,0)
    {\iterbadge{#1}\ \textbf{\color{it#1}Itération #1} --- \itername{#1} \hfill{\color{it#1}\iterdate{#1}}};
  \end{tikzpicture}\par\medskip}

% Encadré d'évolution : \begin{evolution}{V1}{titre} ... \end{evolution}
\newtcolorbox{evolution}[2]{enhanced,breakable,boxrule=0pt,leftrule=3pt,arc=0pt,
  left=6pt,right=6pt,top=4pt,bottom=4pt,colback=it#1!7,colframe=it#1,
  fonttitle=\bfseries\small,coltitle=it#1,attach title to upper={\ },
  title={\iterbadge{#1}\ Évolution #1 : #2.}}

% Légende des itérations, pour les avant-propos
\newcommand{\legendeiterations}{%
\begin{center}\small
\begin{tabular}{lll}
\iterbadge{V0} & \itername{V0} & \iterdate{V0}\\
\iterbadge{V1} & \itername{V1} & \iterdate{V1}\\
\iterbadge{V2} & \itername{V2} & \iterdate{V2}\\
\iterbadge{V3} & \itername{V3} & \iterdate{V3}\\
\end{tabular}
\end{center}}

% Encadré « pistes d'amélioration » : défini ici s'il manque (cours V0 sans cet encadré)
\providecolor{cviolet}{HTML}{7A4FB3}
\ifcsname pistes\endcsname\else
\newtcolorbox{pistes}[1][]{enhanced,breakable,boxrule=0pt,leftrule=2.5pt,arc=0pt,
  left=6pt,right=6pt,top=4pt,bottom=4pt,fonttitle=\bfseries\small,coltitle=cink,attach title to upper={\ },
  colback=cviolet!6,colframe=cviolet,title={Pistes d'amélioration#1.}}
\fi


\title{\textbf{La surface équivalente radar d'un profil}\\[0.4em]\large Un modèle simplifié par optique physique, son implémentation\\ et son intégration dans la chaîne d'optimisation du hackathon quantique}
\author{Notes de travail, projet \emph{HackathonQuantique}}
\date{Version du 4 novembre 2026}

\begin{document}
\pagenumbering{Alph}
\maketitle
\pagenumbering{roman}

\chapter*{Avant-propos}
\addcontentsline{toc}{chapter}{Avant-propos}
\paragraph{Objet.} Le second objectif du hackathon est la surface équivalente radar (SER) du profil. Le modèle définitif sera construit par un autre groupe, à partir d'un maillage facetté et d'un solveur électromagnétique. Pour construire et tester toute la chaîne d'optimisation dès maintenant, on a écrit un modèle \textbf{simplifié mais physique} : l'optique physique (PO) en 2D, monostatique, sur le contour même qu'utilise l'aérodynamique. Ce cours explique ce modèle de la théorie au code, montre ce qu'il dit des profils NACA, et précise comment le vrai solveur viendra le remplacer.

\paragraph{Ce qu'on suppose connu.} Les équations de Maxwell et la notion d'onde plane. Le reste est posé ici : définition de la SER, régimes de diffusion, courants de l'optique physique, intégrale de rayonnement 2D, phase stationnaire.

\paragraph{Fil conducteur.} Définition et ordres de grandeur (ch.~\ref{ch:intro}) $\to$ mécanismes de diffusion haute fréquence et principes de furtivité (ch.~\ref{ch:mecanismes}) $\to$ théorie de l'optique physique (ch.~\ref{ch:po}) $\to$ implémentation et validation (ch.~\ref{ch:code}) $\to$ SER des profils NACA (ch.~\ref{ch:naca}) $\to$ intégration dans la chaîne (ch.~\ref{ch:integration}) $\to$ limites et feuille de route (ch.~\ref{ch:limites}).

\paragraph{Figures et chiffres.} Les courbes à données sont calculées par le code du projet (\texttt{Code/ser/} et \texttt{Code/aero/}) via \texttt{tools/make\_figures.py}. Les chiffres cités sont dans \texttt{figures/chiffres.json}. Chaque chapitre se termine par des exercices (corrigés en annexe~\ref{annexe:corriges}) et par un encadré violet de \textbf{pistes d'amélioration}.

\paragraph{Conventions.} Conducteur électrique parfait. Dépendance temporelle $e^{i\omega t}$ omise. $k=2\pi/\lambda$. Repère du profil : $x$ le long de la corde vers l'aval, $y$ vers l'extrados, origine au bord d'attaque. Corde $c=1$~m, fréquence de référence 10~GHz ($\lambda=3$~cm, $kc\approx209$).

\paragraph{Itérations du projet.} Le projet avance par itérations, chacune avec sa couleur dans tous les cours et dans la fiche projet (dossier \texttt{Fiche\_Projet/} du projet). Chaque chapitre commence par un bandeau qui indique l'itération de son contenu ; un ajout ou une révision faite lors d'une itération ultérieure apparaît dans un encadré « Évolution » de la couleur correspondante.
\legendeiterations

\tableofcontents
\cleardoublepage
\pagenumbering{arabic}

\chapter{La SER : définition et ordres de grandeur}\label{ch:intro}
\iterbanner{V1}

\section{Définition}
Un radar émet une onde, la cible en renvoie une partie, et le radar mesure l'écho. La SER $\sigma$ est la surface fictive qui, interceptant la puissance incidente et la réémettant de façon isotrope, produirait au radar le même écho que la cible :
\[
\sigma=\lim_{R\to\infty}4\pi R^2\,\frac{|\vect E_s|^2}{|\vect E_i|^2}\qquad[\mathrm{m}^2],
\]
où $\vect E_i$ est le champ incident sur la cible et $\vect E_s$ le champ diffusé à la distance $R$. On l'exprime en décibels relatifs à 1~m\textsuperscript{2} : $\sigma_{\mathrm{dBsm}}=10\log_{10}(\sigma/1\,\mathrm{m}^2)$.

\paragraph{Monostatique ou bistatique.} Si l'émetteur et le récepteur sont au même endroit, la SER est \textbf{monostatique} : elle ne dépend que de la direction du radar. C'est le cas étudié dans le hackathon (note d'intention : « radar monostatique, à très haute fréquence, selon un angle d'incidence correspondant à une approche de face »).

\section{Le cas 2D : largeur équivalente radar}
Le profil est une section d'aile : l'objet électromagnétique naturel est un cylindre infini de section le profil. Son écho décroît en $1/\sqrt\rho$ (onde cylindrique) et non en $1/R$. On définit alors la \textbf{largeur équivalente radar} (\emph{echo width}), par unité d'envergure :
\[
\sigma_{2D}=\lim_{\rho\to\infty}2\pi\rho\,\frac{|\vect E_s|^2}{|\vect E_i|^2}\qquad[\mathrm{m}],\qquad
\sigma_{2D,\mathrm{dB}}=10\log_{10}(\sigma_{2D}/1\,\mathrm{m})\quad[\mathrm{dBm}].
\]
Pour un tronçon d'aile d'envergure $b$, vu perpendiculairement à l'envergure, on repasse en 3D par la formule du cylindre fini :
\[
\sigma_{3D}\approx\frac{2b^2}{\lambda}\,\sigma_{2D}=\frac{k\,b^2}{\pi}\,\sigma_{2D}.
\]
Le facteur $b^2/\lambda$ vient de la cohérence de l'écho sur toute l'envergure. Il ne dépend pas du profil : pour \emph{comparer} des profils, $\sigma_{2D}$ suffit, et c'est ce que fait notre modèle.

\section{Les régimes de diffusion}
Tout dépend du rapport entre la taille de l'objet $a$ et la longueur d'onde, mesuré par $ka$ :
\begin{center}
\begin{tabular}{lll}
\toprule
Régime & $ka$ & Comportement \\
\midrule
Rayleigh & $\ll1$ & $\sigma\propto k^4a^6$ en 3D : l'objet est vu comme un dipôle \\
résonance & $\sim1$ à $10$ & ondes rampantes, interférences : il faut résoudre Maxwell (MoM) \\
optique & $\gg1$ & l'écho est formé localement : points spéculaires, arêtes \\
\bottomrule
\end{tabular}
\end{center}
Pour une corde de 1~m, $kc=21$ à 1~GHz et 209 à 10~GHz (bande X). Le profil est donc en régime optique : c'est le sens de « très haute fréquence » dans la note d'intention. Mais attention, le bord d'attaque, lui, a un rayon de 4 à 44~mm selon l'épaisseur, soit $k\,r_{BA}$ de 0,8 à 9 à 10~GHz. Le chapitre~\ref{ch:po} chiffre ce que cela coûte à l'optique physique.

\section{La géométrie du problème}
\begin{figure}[htbp]
\centering
\begin{tikzpicture}[scale=1.15]
\draw[cmuted,dashed] (-5,0) -- (4.5,0) node[right,font=\small] {trajectoire};
\begin{scope}[rotate=8]
\fill[cmuted!25,draw=cink,thick] plot[smooth,tension=0.6] coordinates {(3,0) (2.2,0.16) (1.2,0.28) (0.5,0.27) (0.12,0.15) (0,0) (0.12,-0.1) (0.5,-0.13) (1.2,-0.12) (2.2,-0.05) (3,0)};
\draw[cink,thin] (0,0) -- (3.6,0) node[right,font=\scriptsize] {corde};
\end{scope}
\draw[corange] (2.6,0) arc (0:8:2.6); \node[corange,font=\small] at (2.9,0.17) {$\alpha$};
\fill[cblue] (-4.6,-1.6) circle (2pt) node[below,font=\small,cblue] {radar};
\draw[cblue,-{Stealth[length=2.5mm]},thick] (-4.6,-1.6) -- (-0.15,-0.05);
\draw[cblue,dashed] (-4.6,0) -- (-4.6,-1.6);
\draw[cblue] (-3.2,0) arc (180:199.8:1.4); \node[cblue,font=\small] at (-3.45,-0.3) {$\varepsilon$};
\draw[cgreen!60!black,thick,-{Stealth[length=2mm]}] (-0.2,0.0) -- (-1.4,-0.42) node[left,font=\small] {$\shat$};
\fill[cgreen!12] (-4.9,0) -- (-4.9,1.75) -- (-0.3,0) -- cycle;
\fill[cgreen!12] (-4.9,0) -- (-4.9,-1.75) -- (-0.3,0) -- cycle;
\node[cgreen!50!black,font=\small,align=center] at (-3.6,1.0) {secteur d'approche\\ de face};
\end{tikzpicture}
\caption{Géométrie du problème. Le profil vole à l'incidence $\alpha$ donnée par l'aérodynamique. Le radar est devant l'avion, à l'élévation $\varepsilon$ par rapport à la trajectoire ($\varepsilon<0$ : radar en contrebas). $\shat$ est la direction unitaire du profil vers le radar.}\label{fig:geometrie}
\end{figure}

Dans le repère du profil, le vent relatif vient de la direction $(\cos\alpha,\sin\alpha)$. Un radar à l'élévation $\varepsilon$ est donc vu dans la direction
\[
\shat=-\big(\cos(\alpha-\varepsilon),\ \sin(\alpha-\varepsilon)\big).
\]
L'incidence $\alpha$ n'est \emph{pas} appliquée à la géométrie : c'est la direction radar qui est tournée, comme le prévoit l'interface de la phase~1.

\section{La grandeur optimisée}
La SER varie très vite avec l'angle (lobes, interférences). Un objectif d'optimisation doit être robuste : on moyenne donc $\sigma_{2D}$ \emph{en linéaire} sur un secteur d'élévations, puis on passe en décibels :
\[
\mathrm{SER}=10\log_{10}\Big(\frac{1}{N}\sum_{j=1}^N\sigma_{2D}(\varepsilon_j)\Big),\qquad \varepsilon_j\in[\varepsilon_{\min},\varepsilon_{\max}].
\]
Deux scénarios sont étudiés : S1, conforme à la note d'intention, avec $\varepsilon\in[-20^\circ,20^\circ]$ ; S2, un radar sol vu d'un avion qui l'approche, avec $\varepsilon\in[-60^\circ,0^\circ]$.

\begin{piege}
Moyenner les décibels au lieu des valeurs linéaires sous-estime fortement l'effet des lobes : la moyenne de $-30$~dB et de $0$~dB vaut $-3$~dB en linéaire, mais $-15$~dB en moyenne de décibels. C'est la puissance reçue qui compte, donc la moyenne linéaire.
\end{piege}

\exercices
\begin{exercice}
Le modèle donne $\sigma_{2D}=-13$~dBm pour un profil de corde 1~m à 10~GHz. Estimer la SER 3D d'un tronçon d'aile de 2~m d'envergure, vu perpendiculairement à l'envergure. Pourquoi cette valeur est-elle bien plus grande que celle d'un avion furtif entier vu de face ?
\end{exercice}
\begin{exercice}
Calculer $k\,r_{BA}$ à 10~GHz pour $t=6$~\%, 12~\% et 20~\% ($r_{BA}=1{,}1019\,t^2c$, $c=1$~m). À partir de quelle fréquence le bord d'attaque du profil le plus mince est-il lui aussi en régime optique ($k\,r_{BA}\ge5$) ?
\end{exercice}

\begin{pistes}
\begin{itemize}
\item Fixer avec le groupe SER la corde physique, la fréquence et le secteur : ce sont des paramètres du \emph{problème}, au même titre que le Reynolds côté aérodynamique.
\item Remplacer la moyenne par une statistique plus proche de l'usage opérationnel : percentile 90, ou moyenne pondérée par la probabilité de présence de la menace.
\item Étudier plusieurs fréquences (bandes S, X, Ku) : en 2D et en régime optique, $\sigma_{2D}$ dépend peu de la fréquence, mais $\sigma_{3D}$ croît avec $k$.
\end{itemize}
\end{pistes}

\chapter{Mécanismes de diffusion et principes de furtivité}\label{ch:mecanismes}
\iterbanner{V1}

\section{D'où vient l'écho en régime optique}
En régime optique, l'écho n'est pas « réparti » sur tout l'objet : il est formé en quelques endroits, chacun avec sa loi propre (figure~\ref{fig:mecanismes}).
\begin{itemize}
\item \textbf{Réflexion spéculaire} : là où la normale à la surface pointe vers le radar. Sur une surface courbe de rayon de courbure $\rho$, elle donne $\sigma_{2D}=\pi\rho$, indépendamment de la fréquence (ch.~\ref{ch:po}). Sur une facette plane de largeur $w$ vue de face, $\sigma_{2D}=kw^2$, énorme et très directif.
\item \textbf{Diffraction par les arêtes} : un bord de fuite aigu ou une arête de facette diffracte dans un cône de directions. C'est un écho faible mais présent sous tous les angles.
\item \textbf{Ondes rampantes} : en polarisation TE surtout, une onde contourne l'objet par l'ombre et revient vers le radar. C'est important quand le rayon de courbure est de l'ordre de $\lambda$.
\item \textbf{Ondes progressives, cavités, multi-réflexions} : entrées d'air, jonctions, dièdres. Absentes d'un profil isolé.
\end{itemize}

\begin{figure}[htbp]
\centering
\begin{tikzpicture}[scale=1.3]
\fill[cmuted!25,draw=cink,thick] plot[smooth,tension=0.6] coordinates {(5,0) (3.7,0.27) (2,0.47) (0.8,0.45) (0.2,0.25) (0,0) (0.2,-0.17) (0.8,-0.22) (2,-0.2) (3.7,-0.08) (5,0)};
\draw[cblue,very thick,-{Stealth[length=2.5mm]}] (-2.2,0) -- (-0.1,0);
\draw[corange,very thick,-{Stealth[length=2.5mm]}] (-0.1,0.08) -- (-2.2,0.08);
\node[font=\small,align=center] at (-1.2,0.45) {spéculaire au BA\\ $\sigma_{2D}=\pi r_{BA}$};
\draw[cblue,thick,-{Stealth[length=2mm]}] (1.2,1.6) -- (1.0,0.52);
\draw[corange,thick,-{Stealth[length=2mm]}] (1.05,0.55) -- (1.25,1.62);
\node[font=\small,align=left,anchor=west] at (1.35,1.5) {spéculaire sur l'extrados\\ (radar vu de dessus)};
\draw[cblue,thick,-{Stealth[length=2mm]}] (6.3,0.9) -- (5.05,0.03);
\foreach \a in {150,180,210,240}{\draw[corange,thick,-{Stealth[length=1.6mm]}] (5.05,0) -- ++(\a:0.7);}
\node[font=\small,align=center] at (6.3,-0.6) {diffraction\\ par le bord de fuite};
\draw[cviolet,thick,dashed,-{Stealth[length=2mm]}] (0.15,-0.25) .. controls (1.5,-0.55) and (3.5,-0.4) .. (4.8,-0.1);
\node[cviolet,font=\small] at (2.6,-0.75) {onde rampante (TE)};
\end{tikzpicture}
\caption{Mécanismes de diffusion haute fréquence sur un profil. Notre modèle (optique physique) capture la réflexion spéculaire, et une partie seulement de la diffraction d'arête ; il ignore les ondes rampantes.}\label{fig:mecanismes}
\end{figure}

\section{Les deux philosophies de la furtivité}
La note d'intention rappelle l'histoire. Le F-117 (années 1970-1980) est fait de \textbf{facettes planes} : chaque facette ne renvoie un écho fort que dans la direction de sa normale, choisie hors des directions de menace, et les arêtes sont alignées sur quelques directions. Les codes de l'époque ne savaient calculer que des surfaces planes. Le B-2 utilise des \textbf{surfaces courbes} continues, une fois les codes capables de les traiter : un spéculaire sur une surface courbe vaut $\pi\rho$, faible si $\rho$ est petit, et l'écho reste modéré dans toutes les directions.

Deux règles qualitatives en découlent :
\begin{itemize}
\item éviter les grandes surfaces planes ou peu courbées face à la menace (le spéculaire est proportionnel à $\rho$, ou à $kw^2$) ;
\item garder les arêtes aiguës, et les orienter hors du secteur de menace.
\end{itemize}

\section{Ce que cela dit pour un profil vu de face}
Vu de face, le premier point spéculaire d'un profil est son \textbf{bord d'attaque}. Sa contribution vaut $\pi\rho$ au point où la normale pointe vers le radar : pour un radar dans l'axe, c'est $\pi r_{BA}=1{,}1019\,\pi t^2c$. L'épaisseur, qui fixe ce rayon, sera donc le paramètre dominant. Quand le radar s'écarte de l'axe, le point spéculaire glisse sur l'extrados ou l'intrados, où la courbure est plus faible et $\rho$ plus grand : l'écho croît. La cambrure, qui modifie la courbure de ces deux faces, prend alors de l'importance.

\begin{projet}
Ce raisonnement prédit le résultat principal du chapitre~\ref{ch:naca} : pour une approche de face, la SER d'un profil NACA est essentiellement une fonction de $t$. Comme la traînée croît aussi avec $t$, les deux objectifs poussent dans le même sens, et le compromis se joue sur la cambrure, à épaisseur minimale.
\end{projet}

\exercices
\begin{exercice}
Comparer, à 10~GHz, la largeur équivalente radar d'une facette plane de 10~cm vue de face ($kw^2$) et celle d'un bord d'attaque de rayon 1,6~cm ($\pi\rho$). Combien de décibels les séparent ? Que se passe-t-il si la facette est inclinée de $10^\circ$ ?
\end{exercice}

\begin{pistes}
\begin{itemize}
\item Une paramétrisation qui sépare le rayon de bord d'attaque de l'épaisseur (NACA 4 chiffres modifiée, CST) permettrait de découpler les deux objectifs, et ferait apparaître un vrai compromis.
\item Ajouter un profil « facetté » (bord d'attaque en dièdre) comme cas de comparaison, pour illustrer la philosophie F-117 sur le même problème.
\end{itemize}
\end{pistes}

\chapter{L'optique physique}\label{ch:po}
\iterbanner{V1}

\section{L'approximation}
Sur un conducteur parfait, le champ incident induit une densité de courant surfacique $\vect J_s=\vect n\times\vect H$, où $\vect H$ est le champ magnétique \emph{total} à la paroi et $\vect n$ la normale sortante. Ce courant rayonne le champ diffusé. Le problème exact consiste à trouver $\vect J_s$ en résolvant une équation intégrale sur toute la surface (méthode des moments) : c'est ce que fera le solveur du groupe SER.

L'\textbf{optique physique} remplace ce problème par une approximation locale. En régime optique, chaque élément de surface éclairé se comporte comme un morceau de plan infini tangent, sur lequel le champ total vaut deux fois le champ incident (réflexion totale). D'où :
\[
\vect J_s^{PO}=\begin{cases}2\,\vect n\times\vect H_i & \text{si l'élément est éclairé},\\ \vect 0 & \text{s'il est dans l'ombre}.\end{cases}
\]
Il n'y a plus d'inconnue : on calcule directement le champ rayonné par ce courant.

\section{L'intégrale de rayonnement en 2D}
Prenons la polarisation TM : le champ électrique est parallèle à l'envergure, $\vect E_i=\hat{\vect z}\,e^{ik\,\shat\cdot\vect r}$, pour une onde qui arrive \emph{depuis} la direction $\shat$ du radar. Alors $\vect H_i=-\frac1\eta\,\shat\times\vect E_i$ et le courant PO est porté par $\hat{\vect z}$ :
\[
J_z=\frac{2}{\eta}\,(\vect n\cdot\shat)\,e^{ik\,\shat\cdot\vect r'} .
\]
Le champ rayonné par un courant linéique en 2D fait intervenir la fonction de Hankel :
\[
E_{s,z}(\vect r)=-\frac{k\eta}{4}\int_{\text{contour}}J_z(\vect r')\,H_0^{(2)}\big(k|\vect r-\vect r'|\big)\,\dd\ell'.
\]
Loin de l'objet, dans la direction $\shat$ (retour vers le radar), $H_0^{(2)}(k|\vect r-\vect r'|)\approx\sqrt{\frac{2}{\pi k\rho}}\,e^{-i(k\rho-\pi/4)}\,e^{ik\,\shat\cdot\vect r'}$. En reportant, puis en appliquant la définition $\sigma_{2D}=2\pi\rho|E_s|^2$, tous les facteurs se simplifient :
\[
\boxed{\ \sigma_{2D}=k\,\Big|\int_{\text{éclairé}}(\vect n\cdot\shat)\,e^{2ik\,\shat\cdot\vect r'}\,\dd\ell'\Big|^2\ }
\]
Le facteur 2 dans l'exponentielle est l'aller-retour : un point plus proche du radar de $\delta$ renvoie une onde en avance de phase de $2k\delta$. La polarisation TE donne la même expression en optique physique monostatique 2D ; la différence entre les deux polarisations vient de ce que la PO néglige (ondes rampantes, diffraction).

\section{Une facette plane : formule exacte}
Sur un segment de milieu $\vect r_j$, de tangente $\vect t_j$ et de longueur $L_j$, on paramètre $\vect r'=\vect r_j+\tau\,\vect t_j$ avec $\tau\in[-L_j/2,L_j/2]$ :
\[
\int_{-L_j/2}^{L_j/2}e^{2ik\,\shat\cdot(\vect r_j+\tau\vect t_j)}\,\dd\tau
=L_j\,e^{2ik\,\shat\cdot\vect r_j}\,\operatorname{sinc}\big(kL_j\,\vect t_j\cdot\shat\big),\qquad\operatorname{sinc}x=\frac{\sin x}{x}.
\]
D'où la formule implémentée, \emph{exacte} pour un contour polygonal dans le cadre de la PO :
\[
\sigma_{2D}=k\,\Big|\sum_{j\ \text{éclairées}}(\vect n_j\cdot\shat)\,L_j\,\operatorname{sinc}\big(kL_j\,\vect t_j\cdot\shat\big)\,e^{2ik\,\shat\cdot\vect r_j}\Big|^2 .
\]
Une facette est éclairée si $\vect n_j\cdot\shat>0$. Pour un contour convexe, ce test suffit ; pour une partie concave, une facette peut être tournée vers le radar et cachée par une autre (ombrage), ce que le test ignore.

\paragraph{Cas de la bande.} Une bande de largeur $w$ vue sous l'angle $\theta$ par rapport à sa normale donne $\sigma_{2D}=k\,w^2\cos^2\theta\,\operatorname{sinc}^2(kw\sin\theta)$ : un pic de hauteur $kw^2$ en incidence normale, et des lobes de largeur angulaire $\sim\lambda/(2w)$.

\section{Une surface courbe : la phase stationnaire}
Sur une surface lisse, le terme $e^{2ik\,\shat\cdot\vect r'}$ oscille très vite quand $k$ est grand : les contributions se détruisent, sauf près du point où la phase est stationnaire, c'est-à-dire où $\vect t\cdot\shat=0$. C'est le \textbf{point spéculaire}, où la normale pointe vers le radar. Au voisinage, si $\rho$ est le rayon de courbure et $\ell$ l'abscisse depuis ce point, $\shat\cdot\vect r'\approx\shat\cdot\vect r_0-\frac{\ell^2}{2\rho}$ (la surface s'éloigne du radar), et
\[
\int e^{-ik\ell^2/\rho}\,\dd\ell=\sqrt{\frac{\pi\rho}{ik}}
\quad\Longrightarrow\quad
\boxed{\ \sigma_{2D}=k\cdot\frac{\pi\rho}{k}=\pi\rho\ }
\]
La largeur équivalente d'un point spéculaire ne dépend pas de la fréquence, seulement de la courbure locale. Pour un cylindre de rayon $a$, $\sigma_{2D}=\pi a$, et pour un bord d'attaque vu de face, $\sigma_{2D}=\pi r_{BA}$.

\section{Domaine de validité}
La PO suppose que chaque élément « voit » un plan tangent infini, donc des rayons de courbure grands devant $\lambda$. La solution exacte du cylindre (série de Mie) permet de chiffrer l'erreur commise quand $ka$ n'est pas grand (figure~\ref{fig:mie}).

\begin{figure}[htbp]
\centering
% figure générée par tools/make_figures.py : mie
\begin{tikzpicture}
\begin{axis}[width=0.75\linewidth,height=5.6cm,xmode=log,xlabel={$ka$},ylabel={$10\log_{10}\big(\sigma_{\text{exact}}/\pi a\big)$ (dB)},
  xmin=0.3,xmax=60,ymin=-3,ymax=3,legend style={draw=none,font=\scriptsize,at={(0.5,1.03)},anchor=south},legend columns=2,
  xtick={0.5,1,2,5,10,20,50},xticklabels={0{,}5,1,2,5,10,20,50}]
\addplot[cblue,very thick] coordinates {(0.3,3.0536) (0.32394,2.856) (0.3498,2.6569) (0.37772,2.4577) (0.40786,2.2605) (0.44042,2.0675) (0.47557,1.8817) (0.51352,1.7065) (0.55451,1.5452) (0.59876,1.4011) (0.64655,1.2767) (0.69816,1.1731) (0.75388,1.09) (0.81405,1.0248) (0.87902,0.97274) (0.94917,0.92763) (1.0249,0.88244) (1.1067,0.83076) (1.1951,0.76836) (1.2904,0.69468) (1.3934,0.61376) (1.5046,0.53381) (1.6247,0.46472) (1.7544,0.41364) (1.8944,0.38057) (2.0456,0.35711) (2.2089,0.33102) (2.3852,0.29486) (2.5756,0.25268) (2.7811,0.21693) (3.0031,0.19535) (3.2428,0.18084) (3.5016,0.16075) (3.7811,0.13603) (4.0828,0.11893) (4.4087,0.10878) (4.7606,0.094324) (5.1405,0.080318) (5.5508,0.072615) (5.9938,0.062344) (6.4722,0.054096) (6.9888,0.047764) (7.5465,0.040737) (8.1489,0.035924) (8.7992,0.030787) (9.5015,0.026699) (10.26,0.023272) (11.079,0.019894) (11.963,0.017179) (12.918,0.014897) (13.949,0.012863) (15.062,0.011077) (16.264,0.0095313) (17.562,0.0081944) (18.964,0.0070339) (20.477,0.0060423) (22.112,0.0052059) (23.877,0.0044654) (25.782,0.0038417) (27.84,0.0032969) (30.062,0.0028302) (32.461,0.0024302) (35.052,0.0020867) (37.85,0.0017914) (40.87,0.001537) (44.132,0.0013191) (47.655,0.0011319) (51.458,0.00097117) (55.565,0.00083326) (60,0.00071485)};\addlegendentry{polarisation TM ($E$ parallèle à l'envergure)}
\addplot[corange,very thick] coordinates {(0.3,-7.5467) (0.32394,-6.6262) (0.3498,-5.7244) (0.37772,-4.8455) (0.40786,-3.9951) (0.44042,-3.1796) (0.47557,-2.4071) (0.51352,-1.6867) (0.55451,-1.0293) (0.59876,-0.44725) (0.64655,0.045902) (0.69816,0.43529) (0.75388,0.70492) (0.81405,0.83764) (0.87902,0.81492) (0.94917,0.61698) (1.0249,0.22486) (1.1067,-0.37064) (1.1951,-1.1412) (1.2904,-1.9539) (1.3934,-2.4874) (1.5046,-2.3515) (1.6247,-1.5493) (1.7544,-0.54013) (1.8944,0.23041) (2.0456,0.51733) (2.2089,0.22509) (2.3852,-0.53976) (2.5756,-1.1961) (2.7811,-0.92157) (3.0031,-0.047801) (3.2428,0.36708) (3.5016,-0.10435) (3.7811,-0.73835) (4.0828,-0.29237) (4.4087,0.27116) (4.7606,-0.20208) (5.1405,-0.42137) (5.5508,0.19039) (5.9938,-0.19997) (6.4722,-0.14041) (6.9888,0.073423) (7.5465,-0.22226) (8.1489,0.08561) (8.7992,-0.13815) (9.5015,-0.029424) (10.26,0.050976) (11.079,-0.13519) (11.963,-0.047318) (12.918,0.050635) (13.949,0.046808) (15.062,0.022663) (16.264,0.018581) (17.562,0.029366) (18.964,0.022559) (20.477,-0.027788) (22.112,-0.017596) (23.877,0.0069264) (25.782,-0.0086108) (27.84,-0.020276) (30.062,-0.0053373) (32.461,-0.00306) (35.052,-0.010171) (37.85,-0.0059077) (40.87,0.0006371) (44.132,0.0014111) (47.655,-0.0013833) (51.458,0.0012368) (55.565,-0.001697) (60,0.00086384)};\addlegendentry{polarisation TE}
\addplot[cmuted,dashed] coordinates {(0.3,0) (60,0)};
\draw[cgreen!60!black,thick] (axis cs:0.831,-3) -- (axis cs:0.831,-1.2) node[pos=1,anchor=west,rotate=90,font=\scriptsize] {$t=6$\,\%};
\draw[cgreen!60!black,thick] (axis cs:3.326,-3) -- (axis cs:3.326,-1.2) node[pos=1,anchor=west,rotate=90,font=\scriptsize] {$t=12$\,\%};
\draw[cgreen!60!black,thick] (axis cs:9.238,-3) -- (axis cs:9.238,-1.2) node[pos=1,anchor=west,rotate=90,font=\scriptsize] {$t=20$\,\%};
\end{axis}
\end{tikzpicture}

\caption{Écart entre la SER exacte d'un cylindre conducteur (série de Mie) et la valeur PO $\pi a$, selon $ka$. Traits verticaux : $k\,r_{BA}$ du bord d'attaque à 10~GHz pour trois épaisseurs. Pour $t\ge12$~\%, l'erreur est inférieure à 0,4~dB ; pour $t=6$~\%, elle atteint environ 1~dB.}\label{fig:mie}
\end{figure}

\begin{resultat}
\begin{itemize}
\item PO : courants $2\vect n\times\vect H_i$ sur la partie éclairée, nuls dans l'ombre.
\item En 2D monostatique : $\sigma_{2D}=k|\int(\vect n\cdot\shat)e^{2ik\shat\cdot\vect r'}\dd\ell'|^2$, exacte par facette avec un sinc.
\item Point spéculaire d'une surface courbe : $\sigma_{2D}=\pi\rho$.
\item Validité : $k\rho\gg1$. À 10~GHz, le bord d'attaque des profils minces est à la limite (erreur jusqu'à environ 1~dB).
\end{itemize}
\end{resultat}

\exercices
\begin{exercice}
Retrouver la formule de la bande à partir de la formule par facettes, en considérant une bande comme un rectangle très plat (deux facettes de largeur $w$ dos à dos). Pourquoi une seule des deux facettes contribue-t-elle ?
\end{exercice}
\begin{exercice}
Un cylindre elliptique de demi-axes $a$ (selon $x$) et $b$ (selon $y$) est vu selon l'axe $x$. Le rayon de courbure au sommet vaut $b^2/a$. Que vaut $\sigma_{2D}$ en PO ? Application : un profil elliptique de corde 1~m et d'épaisseur 12~cm, vu de face.
\end{exercice}

\begin{pistes}
\begin{itemize}
\item Ajouter la \textbf{théorie physique de la diffraction} (PTD, Ufimtsev) ou les coefficients de diffraction d'arête (GTD/UTD, Keller) pour représenter le bord de fuite.
\item Distinguer les polarisations TM et TE avec une méthode qui tient compte des ondes rampantes, au moins au bord d'attaque.
\item Calculer la solution exacte par une méthode des moments 2D (équation intégrale du champ électrique, en TM) : quelques centaines d'inconnues suffisent, et cela donnerait une référence pour valider la PO sur les profils eux-mêmes.
\end{itemize}
\end{pistes}

\chapter{Implémentation et validation}\label{ch:code}
\iterbanner{V1}

\section{Organisation du code}
\begin{center}
\begin{tabularx}{\linewidth}{lX}
\toprule
Fichier (\texttt{Code/ser/}) & Contenu \\
\midrule
\code{po2d.py} & facettes, largeur équivalente PO, directions radar, conversions dB et 2D $\to$ 3D \\
\code{model.py} & \code{RadarCase} (scénario), \code{evaluate\_ser} $\to$ \code{SERResult} \\
\bottomrule
\end{tabularx}
\end{center}

\section{Le cœur : \code{echo\_width}}
\begin{codebox}[ : \code{ser/po2d.py}]
\begin{lstlisting}
def facets(X, Y):
    dx, dy = np.diff(X), np.diff(Y)
    L = np.hypot(dx, dy)
    keep = L > 1e-14                       # retire les segments degeneres
    ...
    return xs + dx/2, ys + dy/2, tx, ty, ty, -tx, L    # milieux, tangente, normale sortante, longueur

def echo_width(X, Y, s_dirs, k):
    xc, yc, tx, ty, nx, ny, L = facets(X, Y)
    s = np.atleast_2d(s_dirs)
    ns = s[:, :1] * nx + s[:, 1:] * ny     # n . s   (n_dir x n_fac)
    ts = s[:, :1] * tx + s[:, 1:] * ty     # t . s
    rs = s[:, :1] * xc + s[:, 1:] * yc     # s . r
    lit = ns > 0                           # test d'eclairement
    amp = np.where(lit, ns * L * np.sinc(k * L * ts / np.pi) * np.exp(2j * k * rs), 0.0)
    return k * np.abs(amp.sum(axis=1)) ** 2
\end{lstlisting}
Tout est vectorisé sur les directions et les facettes. \code{np.sinc} de NumPy est le sinus cardinal \emph{normalisé}, $\sin(\pi x)/(\pi x)$ : d'où la division par $\pi$. La normale sortante $(t_y,-t_x)$ suppose le sens de parcours trigonométrique du contour NACA, le même que pour les panneaux.
\end{codebox}

\paragraph{Les directions.} \code{radar\_directions(alpha, eps)} renvoie $\shat=-(\cos(\alpha-\varepsilon),\sin(\alpha-\varepsilon))$ pour un tableau d'élévations. \code{evaluate\_ser} construit le contour NACA (avec \code{n\_half}=200), le met à l'échelle de la corde physique, calcule $\sigma_{2D}$ sur les 81 élévations du secteur et renvoie la moyenne en dBm, la valeur dans l'axe et le maximum. Un appel coûte quelques millisecondes.

\section{Validation}
\begin{figure}[htbp]
\centering
% figure générée par tools/make_figures.py : validation
\begin{tikzpicture}
\begin{groupplot}[group style={group size=2 by 1,horizontal sep=1.7cm},width=0.48\linewidth,height=5.4cm,
  title style={font=\small},tick label style={font=\scriptsize},legend style={draw=none,font=\scriptsize}]
\nextgroupplot[title={Bande de largeur $w=0{,}3$~m, 10~GHz},xlabel={angle depuis la normale (°)},ylabel={$\sigma_{2D}$ (dBm)},ymin=-30,ymax=16,legend pos=north east]
\addplot[cblue,very thick] coordinates {(-20,-21.039) (-19.85,-18.827) (-19.7,-17.21) (-19.55,-16.007) (-19.4,-15.121) (-19.25,-14.497) (-19.1,-14.105) (-18.95,-13.928) (-18.8,-13.963) (-18.65,-14.214) (-18.5,-14.7) (-18.35,-15.454) (-18.2,-16.534) (-18.05,-18.045) (-17.9,-20.182) (-17.75,-23.394) (-17.6,-29.085) (-17.45,-59.021) (-17.3,-29.478) (-17.15,-23.335) (-17,-19.859) (-16.85,-17.506) (-16.7,-15.799) (-16.55,-14.534) (-16.4,-13.603) (-16.25,-12.945) (-16.1,-12.529) (-15.95,-12.335) (-15.8,-12.359) (-15.65,-12.607) (-15.5,-13.097) (-15.35,-13.866) (-15.2,-14.977) (-15.05,-16.541) (-14.9,-18.779) (-14.75,-22.204) (-14.6,-28.574) (-14.45,-46.16) (-14.3,-26.388) (-14.15,-20.849) (-14,-17.582) (-13.85,-15.337) (-13.7,-13.7) (-13.55,-12.487) (-13.4,-11.599) (-13.25,-10.982) (-13.1,-10.606) (-12.95,-10.456) (-12.8,-10.529) (-12.65,-10.837) (-12.5,-11.403) (-12.35,-12.271) (-12.2,-13.521) (-12.05,-15.298) (-11.9,-17.905) (-11.75,-22.137) (-11.6,-31.798) (-11.45,-30.785) (-11.3,-21.489) (-11.15,-17.144) (-11,-14.356) (-10.85,-12.376) (-10.7,-10.914) (-10.55,-9.8292) (-10.4,-9.046) (-10.25,-8.5212) (-10.1,-8.2323) (-9.95,-8.1708) (-9.8,-8.3408) (-9.65,-8.7599) (-9.5,-9.4631) (-9.35,-10.513) (-9.2,-12.024) (-9.05,-14.218) (-8.9,-17.622) (-8.75,-24.063) (-8.6,-39.691) (-8.45,-21.336) (-8.3,-15.83) (-8.15,-12.542) (-8,-10.262) (-7.85,-8.5856) (-7.7,-7.3307) (-7.55,-6.4019) (-7.4,-5.7456) (-7.25,-5.3328) (-7.1,-5.1505) (-6.95,-5.1988) (-6.8,-5.491) (-6.65,-6.0569) (-6.5,-6.9508) (-6.35,-8.2704) (-6.2,-10.203) (-6.05,-13.168) (-5.9,-18.444) (-5.75,-39.097) (-5.6,-19.691) (-5.45,-13.078) (-5.3,-9.3508) (-5.15,-6.7969) (-5,-4.9107) (-4.85,-3.4754) (-4.7,-2.3801) (-4.55,-1.5634) (-4.4,-0.99125) (-4.25,-0.64698) (-4.1,-0.52772) (-3.95,-0.64372) (-3.8,-1.021) (-3.65,-1.7086) (-3.5,-2.7952) (-3.35,-4.4494) (-3.2,-7.0386) (-3.05,-11.62) (-2.9,-25.386) (-2.75,-14.932) (-2.6,-7.2488) (-2.45,-3.0022) (-2.3,-0.027132) (-2.15,2.2629) (-2,4.114) (-1.85,5.6533) (-1.7,6.9551) (-1.55,8.0666) (-1.4,9.0194) (-1.25,9.8358) (-1.1,10.532) (-0.95,11.12) (-0.8,11.609) (-0.65,12.006) (-0.5,12.315) (-0.35,12.541) (-0.2,12.686) (-0.05,12.752) (0.1,12.739) (0.25,12.647) (0.4,12.475) (0.55,12.222) (0.7,11.884) (0.85,11.457) (1,10.936) (1.15,10.313) (1.3,9.5777) (1.45,8.7179) (1.6,7.7151) (1.75,6.5443) (1.9,5.1696) (2.05,3.5364) (2.2,1.557) (2.35,-0.92457) (2.5,-4.2275) (2.65,-9.1967) (2.8,-20.088) (2.95,-17.982) (3.1,-9.7303) (3.25,-6.0322) (3.4,-3.8187) (3.55,-2.3807) (3.7,-1.4406) (3.85,-0.86379) (4,-0.57768) (4.15,-0.54212) (4.3,-0.73685) (4.45,-1.1561) (4.6,-1.8072) (4.75,-2.7119) (4.9,-3.9119) (5.05,-5.4817) (5.2,-7.558) (5.35,-10.421) (5.5,-14.787) (5.65,-23.761) (5.8,-26.36) (5.95,-16.249) (6.1,-12.019) (6.25,-9.4727) (6.4,-7.7744) (6.55,-6.6116) (6.7,-5.8352) (6.85,-5.3651) (7,-5.1566) (7.15,-5.1859) (7.3,-5.4445) (7.45,-5.9363) (7.6,-6.6792) (7.75,-7.7092) (7.9,-9.091) (8.05,-10.942) (8.2,-13.494) (8.35,-17.304) (8.5,-24.369) (8.65,-36.877) (8.8,-21.296) (8.95,-16.293) (9.1,-13.388) (9.25,-11.459) (9.4,-10.119) (9.55,-9.194) (9.7,-8.5908) (9.85,-8.2576) (10,-8.166) (10.15,-8.3032) (10.3,-8.6692) (10.45,-9.2767) (10.6,-10.154) (10.75,-11.354) (10.9,-12.968) (11.05,-15.173) (11.2,-18.354) (11.35,-23.636) (11.5,-39.546) (11.65,-27.23) (11.8,-20.447) (11.95,-16.912) (12.1,-14.633) (12.25,-13.055) (12.4,-11.944) (12.55,-11.183) (12.7,-10.707) (12.85,-10.48) (13,-10.481) (13.15,-10.706) (13.3,-11.16) (13.45,-11.863) (13.6,-12.851) (13.75,-14.192) (13.9,-16.004) (14.05,-18.524) (14.2,-22.323) (14.35,-29.481) (14.5,-40.644) (14.65,-25.855) (14.8,-20.872) (14.95,-17.936) (15.1,-15.959) (15.25,-14.563) (15.4,-13.576) (15.55,-12.905) (15.7,-12.498) (15.85,-12.326) (16,-12.375) (16.15,-12.642) (16.3,-13.137) (16.45,-13.88) (16.6,-14.915) (16.75,-16.312) (16.9,-18.203) (17.05,-20.855) (17.2,-24.932) (17.35,-33.156) (17.5,-38.01) (17.65,-26.711) (17.8,-22.155) (17.95,-19.381) (18.1,-17.484) (18.25,-16.133) (18.4,-15.17) (18.55,-14.51) (18.7,-14.105) (18.85,-13.928) (19,-13.964) (19.15,-14.211) (19.3,-14.678) (19.45,-15.385) (19.6,-16.369) (19.75,-17.697) (19.9,-19.485)};\addlegendentry{code PO}
\addplot[corange,dashed,thick] coordinates {(-20,-21.038) (-19.85,-18.827) (-19.7,-17.21) (-19.55,-16.007) (-19.4,-15.121) (-19.25,-14.497) (-19.1,-14.105) (-18.95,-13.928) (-18.8,-13.963) (-18.65,-14.214) (-18.5,-14.7) (-18.35,-15.454) (-18.2,-16.534) (-18.05,-18.045) (-17.9,-20.183) (-17.75,-23.394) (-17.6,-29.085) (-17.45,-59.027) (-17.3,-29.478) (-17.15,-23.335) (-17,-19.859) (-16.85,-17.505) (-16.7,-15.799) (-16.55,-14.534) (-16.4,-13.603) (-16.25,-12.945) (-16.1,-12.529) (-15.95,-12.335) (-15.8,-12.359) (-15.65,-12.607) (-15.5,-13.097) (-15.35,-13.866) (-15.2,-14.977) (-15.05,-16.541) (-14.9,-18.779) (-14.75,-22.204) (-14.6,-28.574) (-14.45,-46.158) (-14.3,-26.387) (-14.15,-20.849) (-14,-17.582) (-13.85,-15.337) (-13.7,-13.7) (-13.55,-12.487) (-13.4,-11.599) (-13.25,-10.982) (-13.1,-10.606) (-12.95,-10.456) (-12.8,-10.529) (-12.65,-10.837) (-12.5,-11.403) (-12.35,-12.271) (-12.2,-13.521) (-12.05,-15.298) (-11.9,-17.905) (-11.75,-22.137) (-11.6,-31.798) (-11.45,-30.785) (-11.3,-21.489) (-11.15,-17.144) (-11,-14.356) (-10.85,-12.376) (-10.7,-10.914) (-10.55,-9.8292) (-10.4,-9.046) (-10.25,-8.5212) (-10.1,-8.2323) (-9.95,-8.1708) (-9.8,-8.3408) (-9.65,-8.7599) (-9.5,-9.4631) (-9.35,-10.513) (-9.2,-12.024) (-9.05,-14.218) (-8.9,-17.622) (-8.75,-24.063) (-8.6,-39.69) (-8.45,-21.336) (-8.3,-15.83) (-8.15,-12.542) (-8,-10.262) (-7.85,-8.5856) (-7.7,-7.3307) (-7.55,-6.4019) (-7.4,-5.7456) (-7.25,-5.3328) (-7.1,-5.1505) (-6.95,-5.1988) (-6.8,-5.491) (-6.65,-6.0569) (-6.5,-6.9508) (-6.35,-8.2704) (-6.2,-10.203) (-6.05,-13.168) (-5.9,-18.444) (-5.75,-39.098) (-5.6,-19.691) (-5.45,-13.078) (-5.3,-9.3508) (-5.15,-6.7968) (-5,-4.9107) (-4.85,-3.4754) (-4.7,-2.3801) (-4.55,-1.5634) (-4.4,-0.99125) (-4.25,-0.64698) (-4.1,-0.52772) (-3.95,-0.64372) (-3.8,-1.021) (-3.65,-1.7086) (-3.5,-2.7952) (-3.35,-4.4494) (-3.2,-7.0386) (-3.05,-11.62) (-2.9,-25.386) (-2.75,-14.932) (-2.6,-7.2488) (-2.45,-3.0022) (-2.3,-0.027131) (-2.15,2.2629) (-2,4.114) (-1.85,5.6533) (-1.7,6.9551) (-1.55,8.0666) (-1.4,9.0194) (-1.25,9.8358) (-1.1,10.532) (-0.95,11.12) (-0.8,11.609) (-0.65,12.006) (-0.5,12.315) (-0.35,12.541) (-0.2,12.686) (-0.05,12.752) (0.1,12.739) (0.25,12.647) (0.4,12.475) (0.55,12.222) (0.7,11.884) (0.85,11.457) (1,10.936) (1.15,10.313) (1.3,9.5777) (1.45,8.7179) (1.6,7.7151) (1.75,6.5443) (1.9,5.1696) (2.05,3.5364) (2.2,1.557) (2.35,-0.92457) (2.5,-4.2275) (2.65,-9.1967) (2.8,-20.088) (2.95,-17.982) (3.1,-9.7303) (3.25,-6.0322) (3.4,-3.8187) (3.55,-2.3807) (3.7,-1.4406) (3.85,-0.86379) (4,-0.57768) (4.15,-0.54212) (4.3,-0.73685) (4.45,-1.1561) (4.6,-1.8072) (4.75,-2.7119) (4.9,-3.9119) (5.05,-5.4817) (5.2,-7.558) (5.35,-10.421) (5.5,-14.787) (5.65,-23.761) (5.8,-26.36) (5.95,-16.249) (6.1,-12.019) (6.25,-9.4727) (6.4,-7.7744) (6.55,-6.6116) (6.7,-5.8352) (6.85,-5.3651) (7,-5.1566) (7.15,-5.1859) (7.3,-5.4445) (7.45,-5.9363) (7.6,-6.6792) (7.75,-7.7092) (7.9,-9.091) (8.05,-10.942) (8.2,-13.494) (8.35,-17.304) (8.5,-24.369) (8.65,-36.877) (8.8,-21.296) (8.95,-16.293) (9.1,-13.389) (9.25,-11.459) (9.4,-10.119) (9.55,-9.194) (9.7,-8.5908) (9.85,-8.2576) (10,-8.166) (10.15,-8.3032) (10.3,-8.6692) (10.45,-9.2767) (10.6,-10.154) (10.75,-11.354) (10.9,-12.968) (11.05,-15.173) (11.2,-18.354) (11.35,-23.636) (11.5,-39.545) (11.65,-27.23) (11.8,-20.447) (11.95,-16.912) (12.1,-14.633) (12.25,-13.055) (12.4,-11.944) (12.55,-11.183) (12.7,-10.707) (12.85,-10.48) (13,-10.481) (13.15,-10.706) (13.3,-11.16) (13.45,-11.863) (13.6,-12.851) (13.75,-14.192) (13.9,-16.004) (14.05,-18.524) (14.2,-22.323) (14.35,-29.481) (14.5,-40.645) (14.65,-25.855) (14.8,-20.872) (14.95,-17.936) (15.1,-15.959) (15.25,-14.563) (15.4,-13.576) (15.55,-12.905) (15.7,-12.498) (15.85,-12.326) (16,-12.375) (16.15,-12.642) (16.3,-13.137) (16.45,-13.88) (16.6,-14.915) (16.75,-16.312) (16.9,-18.203) (17.05,-20.855) (17.2,-24.932) (17.35,-33.155) (17.5,-38.011) (17.65,-26.711) (17.8,-22.155) (17.95,-19.381) (18.1,-17.484) (18.25,-16.133) (18.4,-15.17) (18.55,-14.51) (18.7,-14.105) (18.85,-13.928) (19,-13.964) (19.15,-14.211) (19.3,-14.678) (19.45,-15.385) (19.6,-16.369) (19.75,-17.697) (19.9,-19.485)};\addlegendentry{formule exacte}
\nextgroupplot[title={Cylindre : $\sigma_{2D}/(\pi a)$},xlabel={nombre de facettes},xmode=log,ymin=0.5,ymax=1.1,legend pos=south east,
  xtick={16,64,256,1024},xticklabels={16,64,256,1024}]
\addplot[cblue,very thick,mark=*,mark size=1.3pt] coordinates {(16,0.90359) (32,1.0022) (64,0.99753) (128,0.99722) (256,0.99718) (512,0.99716) (1024,0.99716)};\addlegendentry{$a=5$~cm ($ka=10$)}
\addplot[corange,very thick,mark=square*,mark size=1.3pt] coordinates {(16,0.0081236) (32,0.96061) (64,0.92764) (128,1.0003) (256,0.99974) (512,0.99966) (1024,0.99965)};\addlegendentry{$a=20$~cm ($ka=42$)}
\addplot[cmuted,dashed] coordinates {(16,1) (1024,1)};
\end{groupplot}
\end{tikzpicture}

\caption{Validation du code. À gauche, bande plane : le code reproduit la formule exacte de la PO à moins de 0,01~dB. À droite, cylindre discrétisé en facettes : le rapport $\sigma_{2D}/(\pi a)$ converge vers 1 (0,997 et 0,9996 avec 1\,024 facettes).}\label{fig:validation}
\end{figure}

Deux tests de référence (figure~\ref{fig:validation}) :
\begin{itemize}
\item \textbf{la bande} vérifie la formule par facette, y compris le sinc et le test d'éclairement ;
\item \textbf{le cylindre} vérifie la limite de phase stationnaire $\pi a$ sur une surface courbe discrétisée. Il faut assez de facettes pour que chacune soit petite devant la zone de Fresnel $\sqrt{\lambda\rho}$ : avec trop peu de facettes, chacune se comporte comme un petit miroir plan, et la SER fluctue.
\end{itemize}

\section{Convergence en facettage, et fréquence}
\begin{figure}[htbp]
\centering
% figure générée par tools/make_figures.py : convergence
\begin{tikzpicture}
\begin{groupplot}[group style={group size=2 by 1,horizontal sep=1.7cm},width=0.48\linewidth,height=5.4cm,
  title style={font=\small},tick label style={font=\scriptsize},legend style={draw=none,font=\scriptsize}]
\nextgroupplot[title={Facettage (secteur S1, 10~GHz)},xlabel={\texttt{n\_half}},ylabel={SER moyenne (dBm)},legend pos=south east]
\addplot[cmuted,very thick,mark=*,mark size=1.3pt] coordinates {(25,-12.832) (50,-12.938) (100,-12.953) (150,-12.941) (200,-12.94) (300,-12.94) (400,-12.94)};\addlegendentry{NACA 0012}
\addplot[cblue,very thick,mark=square*,mark size=1.3pt] coordinates {(25,-12.827) (50,-12.946) (100,-12.954) (150,-12.948) (200,-12.948) (300,-12.948) (400,-12.948)};\addlegendentry{NACA 2412}
\nextgroupplot[title={Fréquence (NACA 2412, $\alpha=2^\circ$)},xlabel={fréquence (GHz)},xmode=log,
  xtick={1,2,5,10,20,40},xticklabels={1,2,5,10,20,40}]
\addplot[corange,very thick,mark=*,mark size=1.1pt] coordinates {(1,-12.898) (1.1356,-12.948) (1.2897,-12.964) (1.4646,-12.924) (1.6633,-13.041) (1.8889,-13.074) (2.1452,-13.05) (2.4362,-13.025) (2.7666,-13.021) (3.1419,-13.022) (3.5681,-12.996) (4.0521,-12.992) (4.6017,-12.983) (5.2259,-12.962) (5.9348,-12.981) (6.7399,-12.97) (7.6541,-12.955) (8.6924,-12.95) (9.8715,-12.943) (11.211,-12.937) (12.731,-12.94) (14.458,-12.934) (16.419,-12.931) (18.647,-12.926) (21.176,-12.925) (24.048,-12.927) (27.31,-12.923) (31.015,-12.928) (35.222,-12.928) (40,-12.92)};
\end{groupplot}
\end{tikzpicture}

\caption{À gauche, SER moyenne du secteur S1 selon le nombre de points du contour : elle est stable à 0,01~dB dès \code{n\_half}=100. À droite, SER du NACA 2412 de 1 à 40~GHz : elle varie de moins de 0,1~dB, comme le prévoit la loi $\pi\rho$, indépendante de la fréquence.}\label{fig:conv-ser}
\end{figure}

\begin{piege}
Le facettage est un paramètre \emph{numérique}. Avec trop peu de facettes, des reflets spéculaires de facettes apparaissent, qui n'existent pas sur la forme lisse. Ici, la distribution en cosinus place des facettes très fines au bord d'attaque, là où se forme l'écho de face, et la SER converge vite. Ce ne sera pas forcément vrai pour un solveur 3D sur un STL : le groupe SER doit refaire cette étude avec son propre code.
\end{piege}

\exercices
\begin{exercice}
Estimer le nombre d'opérations d'un appel à \code{evaluate\_ser} (400 facettes, 81 directions). Le comparer au coût d'une méthode des moments 2D avec 400 inconnues (résolution d'un système dense).
\end{exercice}

\begin{pistes}
\begin{itemize}
\item Ajouter un test de non-régression automatique (bande, cylindre, valeur de référence sur le NACA 2412) dans \code{validation.py}.
\item Tester un cylindre elliptique, dont le rayon de courbure varie, pour valider la loi $\pi\rho(\varepsilon)$ hors de l'axe.
\item Écrire une version 3D de la PO sur les triangles du STL (formule de Gordon pour l'intégrale sur un triangle), pour se rapprocher du solveur prévu par la note d'intention.
\end{itemize}
\end{pistes}

\chapter{La SER des profils NACA}\label{ch:naca}
\iterbanner{V1}

\section{Diagrammes de rayonnement}
\begin{figure}[htbp]
\centering
% figure générée par tools/make_figures.py : diagrammes
\begin{tikzpicture}
\begin{axis}[width=0.9\linewidth,height=6.2cm,xlabel={élévation du radar $\varepsilon$ (°), $\alpha=0$ ; $\varepsilon>0$ : radar au-dessus},
  ylabel={$\sigma_{2D}$ (dBm)},xmin=-90,xmax=90,ymin=-25,ymax=15,xtick={-90,-60,...,90},
  legend style={draw=none,font=\scriptsize,at={(0.5,0.98)},anchor=north},legend columns=3]
\fill[cgreen!10] (axis cs:-20,-25) rectangle (axis cs:20,15);
\node[cgreen!60!black,font=\scriptsize,anchor=south] at (axis cs:0,-25) {secteur S1};
\addplot[cmuted,thick] coordinates {(-90,7.8355) (-89.5,8.7761) (-89,6.9696) (-88.5,7.9815) (-88,6.7157) (-87.5,6.9294) (-87,6.507) (-86.5,5.9502) (-86,6.1972) (-85.5,5.1522) (-85,5.7828) (-84.5,4.4827) (-84,5.2859) (-83.5,3.8726) (-83,4.7358) (-82.5,3.2945) (-82,4.1588) (-81.5,2.7533) (-81,3.5722) (-80.5,2.2631) (-80,2.9778) (-79.5,1.8309) (-79,2.3683) (-78.5,1.4555) (-78,1.7371) (-77.5,1.1278) (-77,1.0828) (-76.5,0.83038) (-76,0.4106) (-75.5,0.54081) (-75,-0.2623) (-74.5,0.23951) (-74,-0.9116) (-73.5,-0.089206) (-73,-1.5121) (-72.5,-0.45518) (-72,-2.0466) (-71.5,-0.86434) (-71,-2.5076) (-70.5,-1.3096) (-70,-2.8982) (-69.5,-1.782) (-69,-3.2341) (-68.5,-2.27) (-68,-3.5357) (-67.5,-2.7613) (-67,-3.8226) (-66.5,-3.2421) (-66,-4.1078) (-65.5,-3.7004) (-65,-4.4018) (-64.5,-4.1281) (-64,-4.7121) (-63.5,-4.5196) (-63,-5.0433) (-62.5,-4.8686) (-62,-5.3965) (-61.5,-5.1731) (-61,-5.7755) (-60.5,-5.4346) (-60,-6.1819) (-59.5,-5.6578) (-59,-6.6134) (-58.5,-5.8502) (-58,-7.0533) (-57.5,-6.0299) (-57,-7.4766) (-56.5,-6.2286) (-56,-7.8332) (-55.5,-6.487) (-55,-8.0564) (-54.5,-6.8482) (-54,-8.0902) (-53.5,-7.3518) (-53,-7.9569) (-52.5,-8.0001) (-52,-7.7659) (-51.5,-8.6938) (-51,-7.6839) (-50.5,-9.1734) (-50,-7.8801) (-49.5,-9.1787) (-49,-8.4549) (-48.5,-8.7992) (-48,-9.2955) (-47.5,-8.4861) (-47,-9.8819) (-46.5,-8.6698) (-46,-9.6795) (-45.5,-9.4599) (-45,-9.127) (-44.5,-10.29) (-44,-9.1153) (-43.5,-10.14) (-43,-9.9455) (-42.5,-9.4968) (-42,-10.714) (-41.5,-9.6724) (-41,-10.214) (-40.5,-10.708) (-40,-9.7687) (-39.5,-10.824) (-39,-10.583) (-38.5,-10.074) (-38,-11.198) (-37.5,-10.573) (-37,-10.396) (-36.5,-11.425) (-36,-10.704) (-35.5,-10.645) (-35,-11.599) (-34.5,-10.943) (-34,-10.807) (-33.5,-11.739) (-33,-11.29) (-32.5,-10.925) (-32,-11.771) (-31.5,-11.739) (-31,-11.116) (-30.5,-11.638) (-30,-12.143) (-29.5,-11.53) (-29,-11.465) (-28.5,-12.183) (-28,-12.147) (-27.5,-11.597) (-27,-11.832) (-26.5,-12.421) (-26,-12.204) (-25.5,-11.77) (-25,-12.008) (-24.5,-12.51) (-24,-12.396) (-23.5,-11.999) (-23,-12.062) (-22.5,-12.508) (-22,-12.677) (-21.5,-12.39) (-21,-12.195) (-20.5,-12.394) (-20,-12.734) (-19.5,-12.801) (-19,-12.572) (-18.5,-12.4) (-18,-12.486) (-17.5,-12.737) (-17,-12.908) (-16.5,-12.87) (-16,-12.722) (-15.5,-12.636) (-15,-12.683) (-14.5,-12.826) (-14,-12.958) (-13.5,-13.009) (-13,-12.968) (-12.5,-12.898) (-12,-12.866) (-11.5,-12.888) (-11,-12.937) (-10.5,-12.987) (-10,-13.031) (-9.5,-13.056) (-9,-13.054) (-8.5,-13.059) (-8,-13.07) (-7.5,-13.076) (-7,-13.093) (-6.5,-13.097) (-6,-13.112) (-5.5,-13.116) (-5,-13.128) (-4.5,-13.134) (-4,-13.138) (-3.5,-13.148) (-3,-13.153) (-2.5,-13.156) (-2,-13.159) (-1.5,-13.162) (-1,-13.161) (-0.5,-13.167) (0,-13.166) (0.5,-13.167) (1,-13.161) (1.5,-13.162) (2,-13.159) (2.5,-13.156) (3,-13.153) (3.5,-13.148) (4,-13.138) (4.5,-13.134) (5,-13.128) (5.5,-13.116) (6,-13.112) (6.5,-13.097) (7,-13.093) (7.5,-13.076) (8,-13.07) (8.5,-13.059) (9,-13.054) (9.5,-13.056) (10,-13.031) (10.5,-12.987) (11,-12.937) (11.5,-12.888) (12,-12.866) (12.5,-12.898) (13,-12.968) (13.5,-13.009) (14,-12.958) (14.5,-12.826) (15,-12.683) (15.5,-12.636) (16,-12.722) (16.5,-12.87) (17,-12.908) (17.5,-12.737) (18,-12.486) (18.5,-12.4) (19,-12.572) (19.5,-12.801) (20,-12.734) (20.5,-12.394) (21,-12.195) (21.5,-12.39) (22,-12.677) (22.5,-12.508) (23,-12.062) (23.5,-11.999) (24,-12.396) (24.5,-12.51) (25,-12.008) (25.5,-11.77) (26,-12.204) (26.5,-12.421) (27,-11.832) (27.5,-11.597) (28,-12.147) (28.5,-12.183) (29,-11.465) (29.5,-11.53) (30,-12.143) (30.5,-11.638) (31,-11.116) (31.5,-11.739) (32,-11.771) (32.5,-10.925) (33,-11.29) (33.5,-11.739) (34,-10.807) (34.5,-10.943) (35,-11.599) (35.5,-10.645) (36,-10.704) (36.5,-11.425) (37,-10.396) (37.5,-10.573) (38,-11.198) (38.5,-10.074) (39,-10.583) (39.5,-10.824) (40,-9.7687) (40.5,-10.708) (41,-10.214) (41.5,-9.6724) (42,-10.714) (42.5,-9.4968) (43,-9.9455) (43.5,-10.14) (44,-9.1153) (44.5,-10.29) (45,-9.127) (45.5,-9.4599) (46,-9.6795) (46.5,-8.6698) (47,-9.8819) (47.5,-8.4861) (48,-9.2955) (48.5,-8.7992) (49,-8.4549) (49.5,-9.1787) (50,-7.8801) (50.5,-9.1734) (51,-7.6839) (51.5,-8.6938) (52,-7.7659) (52.5,-8.0001) (53,-7.9569) (53.5,-7.3518) (54,-8.0902) (54.5,-6.8482) (55,-8.0564) (55.5,-6.487) (56,-7.8332) (56.5,-6.2286) (57,-7.4766) (57.5,-6.0299) (58,-7.0533) (58.5,-5.8502) (59,-6.6134) (59.5,-5.6578) (60,-6.1819) (60.5,-5.4346) (61,-5.7755) (61.5,-5.1731) (62,-5.3965) (62.5,-4.8686) (63,-5.0433) (63.5,-4.5196) (64,-4.7121) (64.5,-4.1281) (65,-4.4018) (65.5,-3.7004) (66,-4.1078) (66.5,-3.2421) (67,-3.8226) (67.5,-2.7613) (68,-3.5357) (68.5,-2.27) (69,-3.2341) (69.5,-1.782) (70,-2.8982) (70.5,-1.3096) (71,-2.5076) (71.5,-0.86434) (72,-2.0466) (72.5,-0.45518) (73,-1.5121) (73.5,-0.089206) (74,-0.9116) (74.5,0.23951) (75,-0.2623) (75.5,0.54081) (76,0.4106) (76.5,0.83038) (77,1.0828) (77.5,1.1278) (78,1.7371) (78.5,1.4555) (79,2.3683) (79.5,1.8309) (80,2.9778) (80.5,2.2631) (81,3.5722) (81.5,2.7533) (82,4.1588) (82.5,3.2945) (83,4.7358) (83.5,3.8726) (84,5.2859) (84.5,4.4827) (85,5.7828) (85.5,5.1522) (86,6.1972) (86.5,5.9502) (87,6.507) (87.5,6.9294) (88,6.7157) (88.5,7.9815) (89,6.9696) (89.5,8.7761) (90,7.8355)};\addlegendentry{NACA 0012}
\addplot[cblue,thick] coordinates {(-90,7.4713) (-89.5,9.023) (-89,6.2348) (-88.5,7.4244) (-88,5.502) (-87.5,6.4111) (-87,4.8112) (-86.5,5.4988) (-86,3.8528) (-85.5,4.6922) (-85,2.8895) (-84.5,4.1189) (-84,2.0238) (-83.5,3.5748) (-83,1.1447) (-82.5,2.996) (-82,0.44681) (-81.5,2.4042) (-81,-0.016172) (-80.5,1.6809) (-80,-0.29238) (-79.5,0.83987) (-79,-0.36938) (-78.5,-0.091537) (-78,-0.41426) (-77.5,-1.1086) (-77,-0.50319) (-76.5,-2.0671) (-76,-0.69363) (-75.5,-2.8749) (-75,-1.0431) (-74.5,-3.4203) (-74,-1.5263) (-73.5,-3.6855) (-73,-2.1522) (-72.5,-3.7743) (-72,-2.875) (-71.5,-3.7688) (-71,-3.6454) (-70.5,-3.7831) (-70,-4.4236) (-69.5,-3.8496) (-69,-5.1384) (-68.5,-3.9915) (-68,-5.7689) (-67.5,-4.2064) (-67,-6.2854) (-66.5,-4.4704) (-66,-6.7049) (-65.5,-4.7789) (-65,-7.0549) (-64.5,-5.1009) (-64,-7.3524) (-63.5,-5.4234) (-63,-7.6356) (-62.5,-5.7361) (-62,-7.9104) (-61.5,-6.028) (-61,-8.1947) (-60.5,-6.307) (-60,-8.4779) (-59.5,-6.5656) (-59,-8.7406) (-58.5,-6.83) (-58,-8.965) (-57.5,-7.1218) (-57,-9.1015) (-56.5,-7.4684) (-56,-9.1105) (-55.5,-7.9103) (-55,-8.9829) (-54.5,-8.4729) (-54,-8.7537) (-53.5,-9.156) (-53,-8.5202) (-52.5,-9.8564) (-52,-8.4084) (-51.5,-10.331) (-51,-8.5642) (-50.5,-10.31) (-50,-9.0946) (-49.5,-9.824) (-49,-9.9695) (-48.5,-9.299) (-48,-10.815) (-47.5,-9.1998) (-47,-10.894) (-46.5,-9.7776) (-46,-10.184) (-45.5,-10.851) (-45,-9.6609) (-44.5,-11.362) (-44,-10.048) (-43.5,-10.628) (-43,-11.205) (-42.5,-10.023) (-42,-11.579) (-41.5,-10.648) (-41,-10.625) (-40.5,-11.82) (-40,-10.458) (-39.5,-11.338) (-39,-11.681) (-38.5,-10.601) (-38,-11.861) (-37.5,-11.575) (-37,-10.864) (-36.5,-12.172) (-36,-11.63) (-35.5,-11.101) (-35,-12.349) (-34.5,-11.838) (-34,-11.266) (-33.5,-12.4) (-33,-12.172) (-32.5,-11.4) (-32,-12.279) (-31.5,-12.568) (-31,-11.628) (-30.5,-11.989) (-30,-12.809) (-29.5,-12.114) (-29,-11.761) (-28.5,-12.59) (-28,-12.776) (-27.5,-12.016) (-27,-12.101) (-26.5,-12.894) (-26,-12.807) (-25.5,-12.167) (-25,-12.299) (-24.5,-12.979) (-24,-12.965) (-23.5,-12.385) (-23,-12.317) (-22.5,-12.852) (-22,-13.17) (-21.5,-12.811) (-21,-12.458) (-20.5,-12.609) (-20,-13.063) (-19.5,-13.216) (-19,-12.925) (-18.5,-12.636) (-18,-12.688) (-17.5,-13.004) (-17,-13.257) (-16.5,-13.212) (-16,-12.974) (-15.5,-12.801) (-15,-12.824) (-14.5,-13.015) (-14,-13.232) (-13.5,-13.308) (-13,-13.225) (-12.5,-13.092) (-12,-12.988) (-11.5,-12.958) (-11,-13.026) (-10.5,-13.126) (-10,-13.216) (-9.5,-13.272) (-9,-13.272) (-8.5,-13.252) (-8,-13.206) (-7.5,-13.174) (-7,-13.146) (-6.5,-13.145) (-6,-13.148) (-5.5,-13.164) (-5,-13.177) (-4.5,-13.184) (-4,-13.186) (-3.5,-13.181) (-3,-13.175) (-2.5,-13.173) (-2,-13.17) (-1.5,-13.164) (-1,-13.157) (-0.5,-13.15) (0,-13.141) (0.5,-13.133) (1,-13.123) (1.5,-13.113) (2,-13.102) (2.5,-13.089) (3,-13.076) (3.5,-13.062) (4,-13.05) (4.5,-13.032) (5,-13.017) (5.5,-13.005) (6,-12.979) (6.5,-12.966) (7,-12.949) (7.5,-12.922) (8,-12.902) (8.5,-12.887) (9,-12.864) (9.5,-12.831) (10,-12.803) (10.5,-12.787) (11,-12.768) (11.5,-12.744) (12,-12.71) (12.5,-12.669) (13,-12.647) (13.5,-12.64) (14,-12.623) (14.5,-12.572) (15,-12.489) (15.5,-12.428) (16,-12.42) (16.5,-12.466) (17,-12.481) (17.5,-12.386) (18,-12.218) (18.5,-12.107) (19,-12.161) (19.5,-12.294) (20,-12.286) (20.5,-12.07) (21,-11.881) (21.5,-11.946) (22,-12.149) (22.5,-12.104) (23,-11.77) (23.5,-11.614) (24,-11.835) (24.5,-11.976) (25,-11.637) (25.5,-11.331) (26,-11.536) (26.5,-11.767) (27,-11.382) (27.5,-11.044) (28,-11.347) (28.5,-11.522) (29,-10.986) (29.5,-10.822) (30,-11.287) (30.5,-11.1) (31,-10.538) (31.5,-10.844) (32,-11.098) (32.5,-10.422) (33,-10.436) (33.5,-10.946) (34,-10.323) (34.5,-10.108) (35,-10.728) (35.5,-10.162) (36,-9.8445) (36.5,-10.512) (37,-9.8946) (37.5,-9.6282) (38,-10.309) (38.5,-9.5266) (39,-9.4944) (39.5,-10.044) (40,-9.0957) (40.5,-9.4805) (41,-9.5865) (41.5,-8.7442) (42,-9.528) (42.5,-8.8932) (43,-8.6693) (43.5,-9.2912) (44,-8.2302) (44.5,-8.8919) (45,-8.4786) (45.5,-8.0719) (46,-8.7669) (46.5,-7.6493) (47,-8.4163) (47.5,-7.7477) (48,-7.6588) (48.5,-8.0123) (49,-7.0585) (49.5,-7.9721) (50,-6.8149) (50.5,-7.486) (51,-6.8506) (51.5,-6.8159) (52,-6.967) (52.5,-6.2172) (53,-6.9589) (53.5,-5.7901) (54,-6.7286) (54.5,-5.5224) (55,-6.3059) (55.5,-5.3586) (56,-5.7873) (56.5,-5.2408) (57,-5.2535) (57.5,-5.1128) (58,-4.7414) (58.5,-4.944) (59,-4.2705) (59.5,-4.7254) (60,-3.8326) (60.5,-4.454) (61,-3.4209) (61.5,-4.1433) (62,-3.0241) (62.5,-3.7962) (63,-2.6326) (63.5,-3.4283) (64,-2.2469) (64.5,-3.0428) (65,-1.8563) (65.5,-2.6404) (66,-1.4662) (66.5,-2.2293) (67,-1.0757) (67.5,-1.8044) (68,-0.68702) (68.5,-1.3737) (69,-0.30671) (69.5,-0.93195) (70,0.073065) (70.5,-0.4802) (71,0.44139) (71.5,-0.030126) (72,0.80266) (72.5,0.42719) (73,1.1645) (73.5,0.87936) (74,1.5211) (74.5,1.3282) (75,1.8935) (75.5,1.7799) (76,2.28) (76.5,2.2157) (77,2.6813) (77.5,2.6468) (78,3.1133) (78.5,3.0581) (79,3.5486) (79.5,3.4109) (80,3.9736) (80.5,3.7113) (81,4.4082) (81.5,3.9681) (82,4.8415) (82.5,4.1916) (83,5.2607) (83.5,4.4615) (84,5.6537) (84.5,4.8717) (85,5.931) (85.5,5.4379) (86,5.9568) (86.5,6.113) (87,5.7766) (87.5,6.7459) (88,5.7936) (88.5,6.9875) (89,6.4512) (89.5,6.646) (90,7.498)};\addlegendentry{NACA 2412}
\addplot[corange,thick] coordinates {(-90,1.9533) (-89.5,7.6714) (-89,2.7701) (-88.5,6.6906) (-88,2.3476) (-87.5,5.5119) (-87,1.2738) (-86.5,4.7008) (-86,0.70287) (-85.5,4.106) (-85,0.1604) (-84.5,3.3105) (-84,-0.2922) (-83.5,2.5235) (-83,-0.30403) (-82.5,1.6805) (-82,-0.26999) (-81.5,0.6343) (-81,-0.14945) (-80.5,-0.4204) (-80,-0.037508) (-79.5,-1.5105) (-79,-0.10967) (-78.5,-2.4948) (-78,-0.28926) (-77.5,-3.1386) (-77,-0.66376) (-76.5,-3.4757) (-76,-1.2159) (-75.5,-3.4637) (-75,-1.8845) (-74.5,-3.3172) (-74,-2.6849) (-73.5,-3.1569) (-73,-3.503) (-72.5,-3.0519) (-72,-4.2993) (-71.5,-3.0777) (-71,-4.9953) (-70.5,-3.1899) (-70,-5.5206) (-69.5,-3.4126) (-69,-5.9101) (-68.5,-3.7072) (-68,-6.1466) (-67.5,-4.0402) (-67,-6.3168) (-66.5,-4.4104) (-66,-6.4496) (-65.5,-4.7657) (-65,-6.5779) (-64.5,-5.1108) (-64,-6.7417) (-63.5,-5.4219) (-63,-6.9274) (-62.5,-5.6873) (-62,-7.1661) (-61.5,-5.9184) (-61,-7.4394) (-60.5,-6.0984) (-60,-7.7463) (-59.5,-6.2553) (-59,-8.0836) (-58.5,-6.3956) (-58,-8.3959) (-57.5,-6.5452) (-57,-8.6552) (-56.5,-6.7551) (-56,-8.7811) (-55.5,-7.0506) (-55,-8.7093) (-54.5,-7.4943) (-54,-8.464) (-53.5,-8.106) (-53,-8.1138) (-52.5,-8.8237) (-52,-7.8322) (-51.5,-9.4627) (-51,-7.7961) (-50.5,-9.6439) (-50,-8.137) (-49.5,-9.2417) (-49,-8.903) (-48.5,-8.6304) (-48,-9.7916) (-47.5,-8.3504) (-47,-10.051) (-46.5,-8.7486) (-46,-9.4061) (-45.5,-9.7391) (-45,-8.7722) (-44.5,-10.355) (-44,-8.9915) (-43.5,-9.7096) (-43,-10.058) (-42.5,-9.0358) (-42,-10.513) (-41.5,-9.5229) (-41,-9.6102) (-40.5,-10.652) (-40,-9.3663) (-39.5,-10.25) (-39,-10.487) (-38.5,-9.4969) (-38,-10.684) (-37.5,-10.36) (-37,-9.7044) (-36.5,-10.907) (-36,-10.369) (-35.5,-9.8619) (-35,-10.988) (-34.5,-10.503) (-34,-9.9507) (-33.5,-10.97) (-33,-10.785) (-32.5,-10.054) (-32,-10.839) (-31.5,-11.145) (-31,-10.276) (-30.5,-10.565) (-30,-11.336) (-29.5,-10.723) (-29,-10.348) (-28.5,-11.076) (-28,-11.293) (-27.5,-10.588) (-27,-10.604) (-26.5,-11.321) (-26,-11.274) (-25.5,-10.657) (-25,-10.723) (-24.5,-11.34) (-24,-11.39) (-23.5,-10.859) (-23,-10.745) (-22.5,-11.203) (-22,-11.513) (-21.5,-11.2) (-21,-10.845) (-20.5,-10.946) (-20,-11.357) (-19.5,-11.529) (-19,-11.273) (-18.5,-10.975) (-18,-10.978) (-17.5,-11.257) (-17,-11.513) (-16.5,-11.489) (-16,-11.245) (-15.5,-11.06) (-15,-11.07) (-14.5,-11.241) (-14,-11.434) (-13.5,-11.518) (-13,-11.443) (-12.5,-11.29) (-12,-11.18) (-11.5,-11.138) (-11,-11.188) (-10.5,-11.288) (-10,-11.373) (-9.5,-11.433) (-9,-11.432) (-8.5,-11.394) (-8,-11.344) (-7.5,-11.283) (-7,-11.253) (-6.5,-11.226) (-6,-11.224) (-5.5,-11.228) (-5,-11.237) (-4.5,-11.245) (-4,-11.251) (-3.5,-11.246) (-3,-11.244) (-2.5,-11.228) (-2,-11.219) (-1.5,-11.207) (-1,-11.194) (-0.5,-11.179) (0,-11.171) (0.5,-11.152) (1,-11.136) (1.5,-11.127) (2,-11.101) (2.5,-11.089) (3,-11.07) (3.5,-11.049) (4,-11.034) (4.5,-11.01) (5,-10.986) (5.5,-10.969) (6,-10.947) (6.5,-10.918) (7,-10.889) (7.5,-10.867) (8,-10.847) (8.5,-10.821) (9,-10.788) (9.5,-10.751) (10,-10.725) (10.5,-10.695) (11,-10.672) (11.5,-10.644) (12,-10.612) (12.5,-10.571) (13,-10.53) (13.5,-10.488) (14,-10.448) (14.5,-10.411) (15,-10.374) (15.5,-10.342) (16,-10.307) (16.5,-10.267) (17,-10.229) (17.5,-10.186) (18,-10.138) (18.5,-10.082) (19,-10.021) (19.5,-9.9858) (20,-9.9629) (20.5,-9.8996) (21,-9.7862) (21.5,-9.7017) (22,-9.7047) (22.5,-9.7218) (23,-9.63) (23.5,-9.4679) (24,-9.4131) (24.5,-9.4868) (25,-9.4643) (25.5,-9.2575) (26,-9.14) (26.5,-9.2443) (27,-9.2555) (27.5,-8.9978) (28,-8.869) (28.5,-9.023) (29,-8.9812) (29.5,-8.6566) (30,-8.6324) (30.5,-8.8169) (31,-8.5725) (31.5,-8.2939) (32,-8.4924) (32.5,-8.4712) (33,-8.0513) (33.5,-8.1228) (34,-8.2776) (34.5,-7.8344) (35,-7.782) (35.5,-8.0324) (36,-7.589) (36.5,-7.4675) (37,-7.7606) (37.5,-7.2927) (38,-7.1902) (38.5,-7.4776) (39,-6.9375) (39.5,-6.9527) (40,-7.1511) (40.5,-6.5537) (41,-6.7773) (41.5,-6.7265) (42,-6.2057) (42.5,-6.6143) (43,-6.1822) (43.5,-6.0042) (44,-6.3136) (44.5,-5.6399) (45,-5.9484) (45.5,-5.7166) (46,-5.3512) (46.5,-5.7528) (47,-5.045) (47.5,-5.3606) (48,-5.0744) (48.5,-4.7835) (49,-5.1107) (49.5,-4.387) (50,-4.8509) (50.5,-4.2646) (51,-4.3437) (51.5,-4.2782) (52,-3.8245) (52.5,-4.2144) (53,-3.457) (53.5,-3.954) (54,-3.2577) (54.5,-3.5186) (55,-3.1605) (55.5,-3.027) (56,-3.0828) (56.5,-2.5732) (57,-2.9379) (57.5,-2.189) (58,-2.6948) (58.5,-1.8892) (59,-2.3618) (59.5,-1.6468) (60,-1.9633) (60.5,-1.4413) (61,-1.5388) (61.5,-1.2465) (62,-1.1073) (62.5,-1.0429) (63,-0.69137) (63.5,-0.82924) (64,-0.29844) (64.5,-0.58942) (65,0.075308) (65.5,-0.33255) (66,0.41606) (66.5,-0.06394) (67,0.73515) (67.5,0.22371) (68,1.0301) (68.5,0.51172) (69,1.2913) (69.5,0.80555) (70,1.5368) (70.5,1.1147) (71,1.758) (71.5,1.4276) (72,1.9575) (72.5,1.7642) (73,2.1552) (73.5,2.1277) (74,2.3329) (74.5,2.4892) (75,2.4856) (75.5,2.8493) (76,2.6271) (76.5,3.1884) (77,2.7454) (77.5,3.4661) (78,2.8711) (78.5,3.6915) (79,3.0911) (79.5,3.8722) (80,3.4412) (80.5,3.9605) (81,3.8599) (81.5,3.9131) (82,4.2608) (82.5,3.8202) (83,4.5948) (83.5,3.9118) (84,4.7824) (84.5,4.332) (85,4.7094) (85.5,4.8942) (86,4.4376) (86.5,5.2509) (87,4.4585) (87.5,5.1663) (88,5.0461) (88.5,4.7297) (89,5.582) (89.5,4.7697) (90,5.5007)};\addlegendentry{NACA 4415}
\end{axis}
\end{tikzpicture}

\caption{Largeur équivalente radar de trois profils à $\alpha=0$, selon l'élévation du radar (10~GHz). La zone verte est le secteur S1. Au centre, l'écho du bord d'attaque ; plus loin, le point spéculaire glisse sur les faces, où la courbure est plus faible, et l'écho croît.}\label{fig:diagrammes}
\end{figure}

La figure~\ref{fig:diagrammes} montre trois traits communs :
\begin{itemize}
\item \textbf{dans l'axe}, une valeur qui ne dépend que du bord d'attaque : $-13{,}2$~dBm pour le 0012 et le 2412 (même épaisseur), $-11{,}2$~dBm pour le 4415, plus épais ;
\item \textbf{hors de l'axe}, une croissance régulière : le point spéculaire s'éloigne du nez, vers des régions moins courbées ;
\item \textbf{une dissymétrie} pour les profils cambrés : l'extrados, plus bombé, et l'intrados, plus plat, ne présentent pas la même courbure au radar.
\end{itemize}

\section{Le bord d'attaque : $\pi r_{BA}$}
Près du bord d'attaque, la loi d'épaisseur NACA se comporte comme $y\approx0{,}2969\cdot5t\sqrt x$, soit $y^2\approx2\,r_{BA}\,x$ avec $r_{BA}=1{,}1019\,t^2$ : une parabole de rayon au sommet $r_{BA}$. Dans l'axe, la loi du point spéculaire donne
\[
\sigma_{2D}^{\text{nez}}\approx\pi\,r_{BA}=1{,}1019\,\pi\,t^2c\qquad\Longrightarrow\qquad
\sigma_{2D,\mathrm{dB}}^{\text{nez}}\approx-13{,}0+20\log_{10}\frac{t}{0{,}12}.
\]
\textbf{Doubler l'épaisseur coûte 6~dB}. La figure~\ref{fig:rayon} vérifie cette loi sur les profils symétriques : l'écart est de quelques dixièmes de décibel, dû aux contributions des autres parties du profil.

\begin{figure}[htbp]
\centering
% figure générée par tools/make_figures.py : rayon
\begin{tikzpicture}
\begin{axis}[width=0.7\linewidth,height=5.6cm,xlabel={épaisseur relative $t$},ylabel={$\sigma_{2D}$ de nez (dBm)},
  legend style={draw=none,font=\scriptsize},legend pos=south east,
  xticklabel style={/pgf/number format/fixed,/pgf/number format/use comma}]
\addplot[cblue,only marks,mark=*,mark size=1.8pt] coordinates {(0.06,-19.187) (0.07,-17.848) (0.08,-16.688) (0.09,-15.665) (0.1,-14.75) (0.11,-13.922) (0.12,-13.166) (0.13,-12.471) (0.14,-11.827) (0.15,-11.228) (0.16,-10.667) (0.17,-10.141) (0.18,-9.6441) (0.19,-9.1745) (0.2,-8.729)};\addlegendentry{PO, profils symétriques}
\addplot[corange,thick,domain=0.05:0.21,samples=50] {10*log10(pi*1.1019*x^2)};\addlegendentry{$\pi\,r_{BA}=1{,}1019\,\pi t^2$}
\end{axis}
\end{tikzpicture}

\caption{SER dans l'axe de profils symétriques d'épaisseur $t$ (10~GHz), comparée à la loi $\pi r_{BA}$. De 6 à 20~\% d'épaisseur, l'écho passe de $-19{,}2$ à $-8{,}7$~dBm.}\label{fig:rayon}
\end{figure}

\paragraph{Hors de l'axe.} Pour une parabole, le rayon de courbure au point où la normale fait l'angle $\beta$ avec l'axe vaut $\rho(\beta)=r_{BA}/\cos^3\beta$. L'écho croît donc comme $1/\cos^3\beta$ : $+1{,}9$~dB à $30^\circ$, $+9$~dB à $60^\circ$. La figure~\ref{fig:speculaire} compare le résultat PO à $\pi\rho$, avec $\rho$ mesuré au point spéculaire sur la géométrie discrète : la loi locale explique l'essentiel du diagramme.

\begin{figure}[htbp]
\centering
% figure générée par tools/make_figures.py : speculaire
\begin{tikzpicture}
\begin{axis}[width=0.75\linewidth,height=5.6cm,xlabel={élévation $\varepsilon$ (°)},ylabel={$\sigma_{2D}$ (dBm)},
  xmin=-60,xmax=60,legend style={draw=none,font=\scriptsize},legend pos=north west]
\addplot[cblue,very thick] coordinates {(-60,-6.1814) (-59.5,-5.6588) (-59,-6.611) (-58.5,-5.8509) (-58,-7.053) (-57.5,-6.0316) (-57,-7.477) (-56.5,-6.2302) (-56,-7.8344) (-55.5,-6.4866) (-55,-8.0557) (-54.5,-6.8476) (-54,-8.0928) (-53.5,-7.3516) (-53,-7.9601) (-52.5,-7.9999) (-52,-7.7677) (-51.5,-8.6911) (-51,-7.6844) (-50.5,-9.1743) (-50,-7.8807) (-49.5,-9.1814) (-49,-8.4548) (-48.5,-8.8008) (-48,-9.2932) (-47.5,-8.4871) (-47,-9.8825) (-46.5,-8.6699) (-46,-9.6822) (-45.5,-9.4587) (-45,-9.1276) (-44.5,-10.289) (-44,-9.1148) (-43.5,-10.142) (-43,-9.944) (-42.5,-9.4982) (-42,-10.713) (-41.5,-9.6704) (-41,-10.215) (-40.5,-10.706) (-40,-9.769) (-39.5,-10.826) (-39,-10.579) (-38.5,-10.074) (-38,-11.197) (-37.5,-10.572) (-37,-10.397) (-36.5,-11.423) (-36,-10.7) (-35.5,-10.645) (-35,-11.6) (-34.5,-10.94) (-34,-10.806) (-33.5,-11.738) (-33,-11.29) (-32.5,-10.925) (-32,-11.771) (-31.5,-11.737) (-31,-11.116) (-30.5,-11.64) (-30,-12.143) (-29.5,-11.528) (-29,-11.467) (-28.5,-12.186) (-28,-12.143) (-27.5,-11.598) (-27,-11.835) (-26.5,-12.42) (-26,-12.203) (-25.5,-11.77) (-25,-12.008) (-24.5,-12.509) (-24,-12.396) (-23.5,-11.996) (-23,-12.06) (-22.5,-12.509) (-22,-12.675) (-21.5,-12.389) (-21,-12.196) (-20.5,-12.394) (-20,-12.736) (-19.5,-12.8) (-19,-12.571) (-18.5,-12.399) (-18,-12.485) (-17.5,-12.736) (-17,-12.908) (-16.5,-12.87) (-16,-12.722) (-15.5,-12.635) (-15,-12.682) (-14.5,-12.824) (-14,-12.959) (-13.5,-13.008) (-13,-12.967) (-12.5,-12.896) (-12,-12.864) (-11.5,-12.889) (-11,-12.938) (-10.5,-12.986) (-10,-13.031) (-9.5,-13.057) (-9,-13.053) (-8.5,-13.058) (-8,-13.069) (-7.5,-13.075) (-7,-13.094) (-6.5,-13.097) (-6,-13.113) (-5.5,-13.116) (-5,-13.129) (-4.5,-13.134) (-4,-13.139) (-3.5,-13.149) (-3,-13.154) (-2.5,-13.155) (-2,-13.158) (-1.5,-13.161) (-1,-13.161) (-0.5,-13.165) (0,-13.167) (0.5,-13.165) (1,-13.161) (1.5,-13.161) (2,-13.158) (2.5,-13.155) (3,-13.154) (3.5,-13.149) (4,-13.139) (4.5,-13.134) (5,-13.129) (5.5,-13.116) (6,-13.113) (6.5,-13.097) (7,-13.094) (7.5,-13.075) (8,-13.069) (8.5,-13.058) (9,-13.053) (9.5,-13.057) (10,-13.031) (10.5,-12.986) (11,-12.938) (11.5,-12.889) (12,-12.864) (12.5,-12.896) (13,-12.967) (13.5,-13.008) (14,-12.959) (14.5,-12.824) (15,-12.682) (15.5,-12.635) (16,-12.722) (16.5,-12.87) (17,-12.908) (17.5,-12.736) (18,-12.485) (18.5,-12.399) (19,-12.571) (19.5,-12.8) (20,-12.736) (20.5,-12.394) (21,-12.196) (21.5,-12.389) (22,-12.675) (22.5,-12.509) (23,-12.06) (23.5,-11.996) (24,-12.396) (24.5,-12.509) (25,-12.008) (25.5,-11.77) (26,-12.203) (26.5,-12.42) (27,-11.835) (27.5,-11.598) (28,-12.143) (28.5,-12.186) (29,-11.467) (29.5,-11.528) (30,-12.143) (30.5,-11.64) (31,-11.116) (31.5,-11.737) (32,-11.771) (32.5,-10.925) (33,-11.29) (33.5,-11.738) (34,-10.806) (34.5,-10.94) (35,-11.6) (35.5,-10.645) (36,-10.7) (36.5,-11.423) (37,-10.397) (37.5,-10.572) (38,-11.197) (38.5,-10.074) (39,-10.579) (39.5,-10.826) (40,-9.769) (40.5,-10.706) (41,-10.215) (41.5,-9.6704) (42,-10.713) (42.5,-9.4982) (43,-9.944) (43.5,-10.142) (44,-9.1148) (44.5,-10.289) (45,-9.1276) (45.5,-9.4587) (46,-9.6822) (46.5,-8.6699) (47,-9.8825) (47.5,-8.4871) (48,-9.2932) (48.5,-8.8008) (49,-8.4548) (49.5,-9.1814) (50,-7.8807) (50.5,-9.1743) (51,-7.6844) (51.5,-8.6911) (52,-7.7677) (52.5,-7.9999) (53,-7.9601) (53.5,-7.3516) (54,-8.0928) (54.5,-6.8476) (55,-8.0557) (55.5,-6.4866) (56,-7.8344) (56.5,-6.2302) (57,-7.477) (57.5,-6.0316) (58,-7.053) (58.5,-5.8509) (59,-6.611) (59.5,-5.6588) (60,-6.1814)};\addlegendentry{PO (NACA 0012)}
\addplot[corange,dashed,very thick] coordinates {(-60,-5.8282) (-59.5,-6.0878) (-59,-6.0878) (-58.5,-6.3507) (-58,-6.3507) (-57.5,-6.6167) (-57,-6.6167) (-56.5,-6.8857) (-56,-6.8857) (-55.5,-7.1575) (-55,-7.1575) (-54.5,-7.4318) (-54,-7.4318) (-53.5,-7.7084) (-53,-7.7084) (-52.5,-7.9869) (-52,-7.9869) (-51.5,-8.2668) (-51,-8.2668) (-50.5,-8.2668) (-50,-8.5478) (-49.5,-8.5478) (-49,-8.8294) (-48.5,-8.8294) (-48,-8.8294) (-47.5,-9.1109) (-47,-9.1109) (-46.5,-9.1109) (-46,-9.3917) (-45.5,-9.3917) (-45,-9.3917) (-44.5,-9.671) (-44,-9.671) (-43.5,-9.671) (-43,-9.9479) (-42.5,-9.9479) (-42,-9.9479) (-41.5,-10.222) (-41,-10.222) (-40.5,-10.222) (-40,-10.491) (-39.5,-10.491) (-39,-10.491) (-38.5,-10.491) (-38,-10.755) (-37.5,-10.755) (-37,-10.755) (-36.5,-11.012) (-36,-11.012) (-35.5,-11.012) (-35,-11.012) (-34.5,-11.26) (-34,-11.26) (-33.5,-11.26) (-33,-11.26) (-32.5,-11.5) (-32,-11.5) (-31.5,-11.5) (-31,-11.5) (-30.5,-11.728) (-30,-11.728) (-29.5,-11.728) (-29,-11.728) (-28.5,-11.943) (-28,-11.943) (-27.5,-11.943) (-27,-11.943) (-26.5,-12.144) (-26,-12.144) (-25.5,-12.144) (-25,-12.144) (-24.5,-12.329) (-24,-12.329) (-23.5,-12.329) (-23,-12.329) (-22.5,-12.329) (-22,-12.497) (-21.5,-12.497) (-21,-12.497) (-20.5,-12.497) (-20,-12.497) (-19.5,-12.645) (-19,-12.645) (-18.5,-12.645) (-18,-12.645) (-17.5,-12.772) (-17,-12.772) (-16.5,-12.772) (-16,-12.772) (-15.5,-12.772) (-15,-12.877) (-14.5,-12.877) (-14,-12.877) (-13.5,-12.877) (-13,-12.877) (-12.5,-12.959) (-12,-12.959) (-11.5,-12.959) (-11,-12.959) (-10.5,-12.959) (-10,-13.017) (-9.5,-13.017) (-9,-13.017) (-8.5,-13.017) (-8,-13.017) (-7.5,-13.05) (-7,-13.05) (-6.5,-13.05) (-6,-13.05) (-5.5,-13.05) (-5,-13.057) (-4.5,-13.057) (-4,-13.057) (-3.5,-13.057) (-3,-13.057) (-2.5,-13.043) (-2,-13.043) (-1.5,-13.043) (-1,-13.043) (-0.5,-13.043) (0,-13.043) (0.5,-13.043) (1,-13.043) (1.5,-13.043) (2,-13.043) (2.5,-13.043) (3,-13.057) (3.5,-13.057) (4,-13.057) (4.5,-13.057) (5,-13.057) (5.5,-13.05) (6,-13.05) (6.5,-13.05) (7,-13.05) (7.5,-13.05) (8,-13.017) (8.5,-13.017) (9,-13.017) (9.5,-13.017) (10,-13.017) (10.5,-12.959) (11,-12.959) (11.5,-12.959) (12,-12.959) (12.5,-12.959) (13,-12.877) (13.5,-12.877) (14,-12.877) (14.5,-12.877) (15,-12.877) (15.5,-12.772) (16,-12.772) (16.5,-12.772) (17,-12.772) (17.5,-12.772) (18,-12.645) (18.5,-12.645) (19,-12.645) (19.5,-12.645) (20,-12.497) (20.5,-12.497) (21,-12.497) (21.5,-12.497) (22,-12.497) (22.5,-12.329) (23,-12.329) (23.5,-12.329) (24,-12.329) (24.5,-12.329) (25,-12.144) (25.5,-12.144) (26,-12.144) (26.5,-12.144) (27,-11.943) (27.5,-11.943) (28,-11.943) (28.5,-11.943) (29,-11.728) (29.5,-11.728) (30,-11.728) (30.5,-11.728) (31,-11.5) (31.5,-11.5) (32,-11.5) (32.5,-11.5) (33,-11.26) (33.5,-11.26) (34,-11.26) (34.5,-11.26) (35,-11.012) (35.5,-11.012) (36,-11.012) (36.5,-11.012) (37,-10.755) (37.5,-10.755) (38,-10.755) (38.5,-10.491) (39,-10.491) (39.5,-10.491) (40,-10.491) (40.5,-10.222) (41,-10.222) (41.5,-10.222) (42,-9.9479) (42.5,-9.9479) (43,-9.9479) (43.5,-9.671) (44,-9.671) (44.5,-9.671) (45,-9.3917) (45.5,-9.3917) (46,-9.3917) (46.5,-9.1109) (47,-9.1109) (47.5,-9.1109) (48,-8.8294) (48.5,-8.8294) (49,-8.8294) (49.5,-8.5478) (50,-8.5478) (50.5,-8.2668) (51,-8.2668) (51.5,-8.2668) (52,-7.9869) (52.5,-7.9869) (53,-7.7084) (53.5,-7.7084) (54,-7.4318) (54.5,-7.4318) (55,-7.1575) (55.5,-7.1575) (56,-6.8857) (56.5,-6.8857) (57,-6.6167) (57.5,-6.6167) (58,-6.3507) (58.5,-6.3507) (59,-6.0878) (59.5,-6.0878) (60,-5.8282)};\addlegendentry{$\pi\rho$ au point spéculaire}
\end{axis}
\end{tikzpicture}

\caption{NACA 0012 : SER PO selon l'élévation, et prédiction $\pi\rho$ au point spéculaire. Les petites oscillations de la courbe PO sont des interférences avec les autres contributions (bord de fuite, faces).}\label{fig:speculaire}
\end{figure}

\section{Sur toute la grille}
\begin{figure}[htbp]
\centering
% figure générée par tools/make_figures.py : cartes
\begin{tikzpicture}
\begin{groupplot}[group style={group size=2 by 1,horizontal sep=3.2cm},width=0.36\linewidth,height=5.6cm,
  xlabel={cambrure $m$ (\%)},ylabel={épaisseur $t$ (\%)},xtick={0,...,7},ytick={6,8,...,20},
  enlargelimits=false,axis on top,title style={font=\small},tick label style={font=\scriptsize},
  colormap={bl}{color=(white) color=(cblue!45) color=(cblue!85!black)}]
\nextgroupplot[title={S1 : face, $\varepsilon\in[-20^\circ,20^\circ]$},colorbar,colorbar style={ylabel={SER (dBm)},width=2.5mm,tick label style={font=\scriptsize}}]
\addplot[matrix plot*,mesh/cols=8,point meta=explicit] coordinates {(0,6) [-18.914] (1,6) [-18.927] (2,6) [-18.901] (3,6) [-18.852] (4,6) [-18.781] (5,6) [-18.687] (6,6) [-18.577] (7,6) [-18.450] (0,8) [-16.427] (1,8) [-16.437] (2,8) [-16.424] (3,8) [-16.394] (4,8) [-16.344] (5,8) [-16.269] (6,8) [-16.166] (7,8) [-16.038] (0,10) [-14.507] (1,10) [-14.517] (2,10) [-14.510] (3,10) [-14.479] (4,10) [-14.429] (5,10) [-14.365] (6,10) [-14.287] (7,10) [-14.192] (0,12) [-12.940] (1,12) [-12.953] (2,12) [-12.948] (3,12) [-12.923] (4,12) [-12.886] (5,12) [-12.831] (6,12) [-12.757] (7,12) [-12.667] (0,14) [-11.622] (1,14) [-11.635] (2,14) [-11.633] (3,14) [-11.617] (4,14) [-11.585] (5,14) [-11.539] (6,14) [-11.481] (7,14) [-11.410] (0,16) [-10.486] (1,16) [-10.497] (2,16) [-10.499] (3,16) [-10.490] (4,16) [-10.467] (5,16) [-10.433] (6,16) [-10.386] (7,16) [-10.324] (0,18) [-9.489] (1,18) [-9.499] (2,18) [-9.503] (3,18) [-9.501] (4,18) [-9.489] (5,18) [-9.465] (6,18) [-9.428] (7,18) [-9.379] (0,20) [-8.601] (1,20) [-8.611] (2,20) [-8.618] (3,20) [-8.620] (4,20) [-8.618] (5,20) [-8.605] (6,20) [-8.580] (7,20) [-8.544]};
\nextgroupplot[title={S2 : radar en contrebas, $\varepsilon\in[-60^\circ,0^\circ]$},colorbar,colorbar style={ylabel={SER (dBm)},width=2.5mm,tick label style={font=\scriptsize}}]
\addplot[matrix plot*,mesh/cols=8,point meta=explicit] coordinates {(0,6) [-15.420] (1,6) [-15.951] (2,6) [-16.406] (3,6) [-16.790] (4,6) [-17.111] (5,6) [-17.373] (6,6) [-17.578] (7,6) [-17.727] (0,8) [-12.973] (1,8) [-13.502] (2,8) [-13.969] (3,8) [-14.377] (4,8) [-14.732] (5,8) [-15.036] (6,8) [-15.292] (7,8) [-15.502] (0,10) [-11.391] (1,10) [-11.881] (2,10) [-12.314] (3,10) [-12.696] (4,10) [-13.030] (5,10) [-13.321] (6,10) [-13.571] (7,10) [-13.781] (0,12) [-10.029] (1,12) [-10.496] (2,12) [-10.908] (3,12) [-11.270] (4,12) [-11.586] (5,12) [-11.862] (6,12) [-12.102] (7,12) [-12.310] (0,14) [-8.828] (1,14) [-9.267] (2,14) [-9.670] (3,14) [-10.030] (4,14) [-10.347] (5,14) [-10.624] (6,14) [-10.865] (7,14) [-11.074] (0,16) [-7.906] (1,16) [-8.284] (2,16) [-8.642] (3,16) [-8.977] (4,16) [-9.282] (5,16) [-9.554] (6,16) [-9.792] (7,16) [-9.998] (0,18) [-7.087] (1,18) [-7.424] (2,18) [-7.735] (3,18) [-8.032] (4,18) [-8.315] (5,18) [-8.581] (6,18) [-8.822] (7,18) [-9.036] (0,20) [-6.379] (1,20) [-6.694] (2,20) [-6.980] (3,20) [-7.238] (4,20) [-7.482] (5,20) [-7.719] (6,20) [-7.948] (7,20) [-8.161]};
\end{groupplot}
\end{tikzpicture}

\caption{SER moyenne (dBm) sur la tranche $p=0{,}40$ de la grille, pour les deux scénarios (incidence $2^\circ$). Dans les deux cas, l'épaisseur domine ; dans S2, où le radar voit surtout l'intrados, la cambrure compte davantage.}\label{fig:cartes}
\end{figure}

\begin{center}
\begin{tabular}{lcc}
\toprule
 & S1 : $\varepsilon\in[-20^\circ,20^\circ]$ & S2 : $\varepsilon\in[-60^\circ,0^\circ]$ \\
\midrule
plage de SER sur les 512 profils & $-18{,}9$ à $-8{,}1$~dBm & $-18{,}8$ à $-6{,}4$~dBm \\
corrélation avec $t$ & 0,98 & 0,95 \\
corrélation avec $m$ & 0,04 & $-0{,}24$ \\
\bottomrule
\end{tabular}
\end{center}

\paragraph{Lecture.} Dans le scénario de la note d'intention, la SER est presque une fonction de la seule épaisseur. La cambrure joue un rôle secondaire, qui devient sensible quand le radar regarde l'intrados (S2) : une cambrure plus forte rend l'intrados plus plat sous le nez, ou même concave à l'arrière, et modifie la courbure présentée au radar.

\begin{projet}
Conséquence directe pour l'optimisation : la traînée et la SER croissent toutes deux avec l'épaisseur. Les deux objectifs poussent donc $t$ vers sa borne $t_{\min}=8$~\%, et le front de Pareto se réduit à des profils de cette épaisseur, qui ne diffèrent que par la cambrure. Dans S1, la SER ne varie que de 0,13~dB le long du front ; dans S2, de 1,7~dB. C'est ce qu'analyse le cours sur l'optimisation.
\end{projet}

\exercices
\begin{exercice}
De combien de décibels la SER de nez diminue-t-elle quand on passe de $t=12$~\% à $t=8$~\% ? Et quelle est, d'après le modèle aérodynamique, la variation de traînée correspondante pour un NACA symétrique à $\alpha=2^\circ$ (ordre de grandeur, voir le cours aéro) ?
\end{exercice}
\begin{exercice}
Vérifier la formule $\rho(\beta)=r/\cos^3\beta$ pour la parabole $y^2=2rx$, en calculant la courbure au point où la tangente fait l'angle $\frac\pi2-\beta$ avec l'axe des $x$.
\end{exercice}

\begin{pistes}
\begin{itemize}
\item Étudier la sensibilité du front à la largeur du secteur et à son centre (menace aérienne ou au sol) : c'est le paramètre qui pilote le poids de la cambrure.
\item Ajouter une variable de conception qui agit sur la SER sans trop agir sur la traînée (rayon de bord d'attaque séparé, bord d'attaque affûté) pour obtenir un vrai compromis.
\end{itemize}
\end{pistes}

\chapter{Intégration dans la chaîne d'optimisation}\label{ch:integration}
\iterbanner{V1}

\section{Ce qui circule entre les phases}
\begin{figure}[htbp]
\centering
\begin{tikzpicture}[node distance=0.5cm and 0.9cm]
\node[flow] (x) {profil $(i_m,i_p,i_t)$\\ 9 bits};
\node[flow,right=of x] (aero) {\textbf{Aéro} (phase 1)\\ \code{evaluate}};
\node[flow,right=of aero,fill=cblue!6] (ser) {\textbf{SER} (phase 2)\\ \code{evaluate\_ser}};
\node[flow,right=of ser] (obj) {objectifs\\ $(C_D,\ \mathrm{SER})$\\ + contraintes};
\draw[arr] (x) -- (aero);
\draw[arr] (aero) -- node[above,font=\scriptsize] {$\alpha$} (ser);
\draw[arr] (ser) -- (obj);
\draw[arr] (aero.south) |- ++(0,-0.8) -| node[pos=0.25,below,font=\scriptsize] {$C_D$, $C_L$, validité} (obj.south);
\node[flow,above=1.1cm of ser,xshift=1.6cm,fill=corange!8] (geo) {géométrie : mêmes points NACA\\ (ou STL de \code{geometry\_export})};
\draw[arr,corange] (geo) -- (ser);
\node[flow,above=1.1cm of aero,xshift=-1.6cm,fill=cgreen!8] (sc) {scénario : $\Re$, $\alpha$ ou $C_L$,\\ $f$, corde, secteur};
\draw[arr,cgreen!60!black] (sc) -- (aero);
\draw[arr,cgreen!60!black] (sc.east) -- (geo.west);
\end{tikzpicture}
\caption{La SER dans la chaîne. Elle reçoit la géométrie et l'attitude $\alpha$ de l'aérodynamique, et un scénario radar partagé. Les optimiseurs ne voient que les deux objectifs et les contraintes.}\label{fig:flux}
\end{figure}

\paragraph{Le couplage par l'attitude.} À incidence imposée (scénario de la note d'intention), tous les profils volent à $2^\circ$ et la SER ne dépend que de la forme. À portance imposée, chaque profil a sa propre incidence (de $-5{,}5^\circ$ à $+3{,}5^\circ$ sur la grille) : l'attitude vue par le radar dépend alors de l'aérodynamique, et c'est un vrai couplage multiphysique. Le code gère les deux, car la SER reçoit toujours l'$\alpha$ calculé par l'aéro.

\section{Le code d'intégration}
\begin{codebox}[ : \code{optim/problem.py}]
\begin{lstlisting}
def simulate(idx, space, sc):
    m, p, t = space.values(idx)
    a = evaluate(m, p, t, alpha=fc.alpha, Re=fc.Re)      # ou evaluate_at_cl en mode portance
    s = evaluate_ser(m, p, t, alpha=a.alpha - offset, case=rc)   # attitude venue de l'aero
    g = [fc.t_min - t, fc.CL_min - a.CL, 0.0 if a.valid else 1.0]  # contraintes g <= 0
    return dict(..., CD_counts=a.CD * 1e4, SER_db=s.ser_db, feasible=all(gi <= 0 for gi in g))
\end{lstlisting}
\code{offset} recentre le secteur d'élévations pour un secteur asymétrique comme S2 : décaler toutes les élévations de $\varepsilon_0$ revient à remplacer $\alpha$ par $\alpha-\varepsilon_0$ dans $\shat$. \code{build\_table} évalue une fois les 512 profils et les enregistre dans \code{resultats/base\_complete\_<scénario>.csv}.
\end{codebox}

\section{Le contrat avec le groupe SER}
Pour que le vrai solveur remplace ce modèle sans rien changer ailleurs, il doit respecter l'interface suivante :
\begin{center}
\renewcommand{\arraystretch}{1.2}
\begin{tabularx}{\linewidth}{lX}
\toprule
Élément & Convention \\
\midrule
Entrée géométrique & STL de \code{geometry\_export} (mêmes points que l'aéro), repère corps : $x$ corde vers l'aval, $y$ extrados, $z$ envergure, origine au bord d'attaque \\
Entrée d'attitude & $\alpha$ (degrés), lu dans le JSON ; la direction radar est tournée, pas la géométrie \\
Scénario & fréquence, corde physique, polarisation, secteur d'élévations, à figer ensemble \\
Sortie & une valeur scalaire par profil : moyenne linéaire sur le secteur, en dB (dBsm en 3D, dBm en 2D) \\
Signature Python & \code{evaluate\_ser(m, p, t, alpha, case) -> objet avec .ser\_db} \\
Coût & compté par le même compteur que l'aéro : une « simulation » = aéro + SER \\
\bottomrule
\end{tabularx}
\end{center}

\paragraph{Procédure de remplacement.}
\begin{enumerate}
\item Écrire un adaptateur \code{evaluate\_ser} qui exporte le STL, appelle le solveur et renvoie \code{ser\_db}.
\item Comparer au modèle PO sur une dizaine de profils : on attend un accord de tendance (rôle de $t$), pas forcément de niveau.
\item Relancer \code{build\_table} pour les deux scénarios : tous les optimiseurs et toutes les figures suivent automatiquement.
\end{enumerate}

\begin{piege}
En 3D, la SER d'un tronçon d'aile extrudé contient l'écho des \emph{flancs} (les deux extrémités du tronçon), qui n'a rien à voir avec le profil. Il faut soit les retirer (\code{extrude(..., caps=False)}), soit les ignorer en ne regardant que les directions perpendiculaires à l'envergure, où ils sont vus en incidence rasante.
\end{piege}

\section{Ce que la SER impose à l'optimisation}
La SER moyenne sur le secteur est une fonction lisse de $(m,p,t)$ dans notre modèle : un modèle quadratique la représente bien, mieux que la traînée. Avec un vrai solveur, deux sources d'irrégularité apparaîtront : le passage d'une facette de l'ombre à la lumière, et les interférences entre contributions, sensibles à la fréquence. C'est l'argument de la note d'intention pour NSGA-II, qui n'utilise pas de gradient. Pour la boucle QUBO, une SER irrégulière dégraderait le modèle quadratique : il faudra le vérifier dès que les vraies valeurs seront disponibles.

\exercices
\begin{exercice}
En mode portance imposée, un profil très cambré vole à $\alpha=-4^\circ$. Dans le scénario S1, le radar le voit-il plutôt par l'extrados ou par l'intrados, par rapport à un profil qui vole à $+2^\circ$ ? Quel effet peut-on attendre sur la SER ?
\end{exercice}

\begin{pistes}
\begin{itemize}
\item Écrire dès maintenant l'adaptateur pour le format d'entrée du solveur choisi par le groupe SER, et le tester avec le modèle PO en guise de solveur factice.
\item Mettre en cache les résultats SER par profil et par scénario (comme \code{AeroSimulator}), pour que les appels au vrai solveur, plus coûteux, ne soient jamais répétés.
\item Prévoir un mode « multi-fidélité » : PO pour explorer, solveur complet pour confirmer les profils du front.
\end{itemize}
\end{pistes}

\chapter{Limites du modèle et feuille de route}\label{ch:limites}
\iterbanner{V1}

\section{Ce que le modèle ne fait pas}
\begin{center}
\renewcommand{\arraystretch}{1.2}
\begin{tabularx}{\linewidth}{lXl}
\toprule
Limite & Effet attendu & Gravité ici \\
\midrule
2D (section infinie) & pas d'effet d'envergure, de flèche ni d'extrémité d'aile & faible pour comparer des profils \\
Optique physique & bord d'attaque des profils minces à la limite ($kr\approx0{,}8$) : environ 1~dB d'erreur & moyenne \\
Pas de diffraction d'arête & bord de fuite sous-estimé, surtout hors des points spéculaires & faible de face \\
Pas d'ondes rampantes & polarisations TM et TE confondues & moyenne pour les profils minces \\
Ombrage par simple test $\vect n\cdot\shat>0$ & faux pour des parties concaves (intrados très cambré) & faible \\
Conducteur parfait & pas de matériaux absorbants & hors sujet ici \\
Secteur moyenné & une seule grandeur scalaire & choix de modélisation \\
\bottomrule
\end{tabularx}
\end{center}

\section{Feuille de route, par rapport utilité / effort}
\begin{enumerate}
\item \textbf{Méthode des moments 2D en polarisation TM} (EFIE) : quelques heures d'implémentation, quelques centaines d'inconnues. Elle donnerait une référence exacte pour valider la PO sur les vrais profils, et lèverait le doute sur le bord d'attaque des profils minces.
\item \textbf{PO 3D sur le STL}, avec la formule fermée de l'intégrale sur un triangle (Gordon, 1975) : c'est le modèle le plus proche de ce que la note d'intention prévoit, et il permet de tester la chaîne STL $\to$ SER.
\item \textbf{Diffraction d'arête} (PTD ou UTD) pour le bord de fuite.
\item \textbf{Étude de sensibilité} à la fréquence, à la polarisation et au secteur, avec le front de Pareto comme sortie.
\end{enumerate}

\begin{projet}
Le message pour l'équipe : le modèle simplifié suffit pour construire, déboguer et comparer les méthodes d'optimisation. Il prédit que, pour une approche de face, la SER d'un profil NACA est dominée par le rayon de bord d'attaque, donc par l'épaisseur. Si le vrai solveur confirme cette tendance, le front de Pareto restera concentré à épaisseur minimale, et c'est la paramétrisation, plus que le solveur, qu'il faudra enrichir pour obtenir un compromis intéressant.
\end{projet}

\exercices
\begin{exercice}
Pourquoi une méthode des moments 2D en TM est-elle beaucoup plus simple à écrire qu'en TE, et pourquoi a-t-on besoin d'environ 10 inconnues par longueur d'onde ? Combien d'inconnues faut-il pour un profil de corde 1~m à 10~GHz ?
\end{exercice}

\begin{pistes}
\begin{itemize}
\item Cette feuille de route est elle-même la liste des pistes : commencer par la méthode des moments 2D, qui sert de validation à tout le reste.
\item Documenter l'écart PO/MoM sur la grille sous forme de corrélation de rang, comme pour la traînée : pour l'optimisation, c'est le classement qui compte.
\end{itemize}
\end{pistes}

\chapter{V2 : le rayon de bord d'attaque comme variable}\label{ch:v2}
\iterbanner{V2}

\begin{evolution}{V2}{pourquoi ce chapitre}En V1, la SER de face d'un profil NACA dépendait presque uniquement de l'épaisseur, par le rayon de bord d'attaque $r_{BA}\propto t^2$ (chapitre~\ref{ch:naca}). En V2, la géométrie devient la série NACA modifiée (cours aéro, chapitre V2), où l'indice $I$ fixe le rayon indépendamment de l'épaisseur : $r_{BA}=1{,}1019\,(tI/6)^2$. Le modèle SER lui-même ne change pas.\end{evolution}

\section{La SER de nez en fonction de $I$}
\begin{figure}[htbp]
\centering
% figure générée par tools/make_figures_v2.py : v2_nez
\begin{tikzpicture}
\begin{axis}[width=0.72\linewidth,height=5.8cm,xlabel={indice de rayon de bord d'attaque $I$},ylabel={$\sigma_{2D}$ de nez (dBm)},
  xmin=1.5,xmax=9,legend style={draw=none,font=\scriptsize},legend pos=south east]
\addplot[corange,very thick] coordinates {(1.5,-25.683) (1.75,-24.879) (2,-24.134) (2.25,-23.442) (2.5,-22.796) (2.75,-22.191) (3,-21.623) (3.25,-21.088) (3.5,-20.583) (3.75,-20.103) (4,-19.648) (4.25,-19.215) (4.5,-18.802) (4.75,-18.406) (5,-18.028) (5.25,-17.665) (5.5,-17.316) (5.75,-16.98) (6,-16.657) (6.25,-16.345) (6.5,-16.043) (6.75,-15.752) (7,-15.47) (7.25,-15.197) (7.5,-14.931) (7.75,-14.674) (8,-14.424) (8.25,-14.181) (8.5,-13.944) (8.75,-13.714) (9,-13.489)};\addlegendentry{PO, $t=8$~\%}
\addplot[corange,dashed,forget plot] coordinates {(1.5,-28.586) (1.75,-27.248) (2,-26.088) (2.25,-25.065) (2.5,-24.15) (2.75,-23.322) (3,-22.566) (3.25,-21.871) (3.5,-21.227) (3.75,-20.628) (4,-20.067) (4.25,-19.541) (4.5,-19.044) (4.75,-18.574) (5,-18.129) (5.25,-17.705) (5.5,-17.301) (5.75,-16.915) (6,-16.545) (6.25,-16.191) (6.5,-15.85) (6.75,-15.522) (7,-15.206) (7.25,-14.902) (7.5,-14.607) (7.75,-14.322) (8,-14.047) (8.25,-13.779) (8.5,-13.52) (8.75,-13.268) (9,-13.023)};
\addplot[cblue,very thick] coordinates {(1.5,-22.161) (1.75,-21.357) (2,-20.612) (2.25,-19.92) (2.5,-19.274) (2.75,-18.67) (3,-18.102) (3.25,-17.566) (3.5,-17.061) (3.75,-16.582) (4,-16.126) (4.25,-15.693) (4.5,-15.28) (4.75,-14.885) (5,-14.506) (5.25,-14.143) (5.5,-13.794) (5.75,-13.458) (6,-13.135) (6.25,-12.823) (6.5,-12.522) (6.75,-12.23) (7,-11.948) (7.25,-11.675) (7.5,-11.41) (7.75,-11.152) (8,-10.902) (8.25,-10.659) (8.5,-10.422) (8.75,-10.192) (9,-9.9676)};\addlegendentry{PO, $t=12$~\%}
\addplot[cblue,dashed,forget plot] coordinates {(1.5,-25.065) (1.75,-23.726) (2,-22.566) (2.25,-21.543) (2.5,-20.628) (2.75,-19.8) (3,-19.044) (3.25,-18.349) (3.5,-17.705) (3.75,-17.106) (4,-16.545) (4.25,-16.019) (4.5,-15.522) (4.75,-15.053) (5,-14.607) (5.25,-14.183) (5.5,-13.779) (5.75,-13.393) (6,-13.023) (6.25,-12.669) (6.5,-12.328) (6.75,-12) (7,-11.685) (7.25,-11.38) (7.5,-11.085) (7.75,-10.8) (8,-10.525) (8.25,-10.257) (8.5,-9.9981) (8.75,-9.7463) (9,-9.5016)};
\addplot[cgreen,very thick] coordinates {(1.5,-19.662) (1.75,-18.858) (2,-18.114) (2.25,-17.421) (2.5,-16.776) (2.75,-16.171) (3,-15.603) (3.25,-15.068) (3.5,-14.562) (3.75,-14.083) (4,-13.628) (4.25,-13.194) (4.5,-12.781) (4.75,-12.386) (5,-12.007) (5.25,-11.644) (5.5,-11.295) (5.75,-10.96) (6,-10.636) (6.25,-10.324) (6.5,-10.023) (6.75,-9.7315) (7,-9.4493) (7.25,-9.176) (7.5,-8.9108) (7.75,-8.6534) (8,-8.4034) (8.25,-8.1603) (8.5,-7.9237) (8.75,-7.6933) (9,-7.4688)};\addlegendentry{PO, $t=16$~\%}
\addplot[cgreen,dashed,forget plot] coordinates {(1.5,-22.566) (1.75,-21.227) (2,-20.067) (2.25,-19.044) (2.5,-18.129) (2.75,-17.301) (3,-16.545) (3.25,-15.85) (3.5,-15.206) (3.75,-14.607) (4,-14.047) (4.25,-13.52) (4.5,-13.023) (4.75,-12.554) (5,-12.108) (5.25,-11.685) (5.5,-11.28) (5.75,-10.894) (6,-10.525) (6.25,-10.17) (6.5,-9.8294) (6.75,-9.5016) (7,-9.1857) (7.25,-8.8809) (7.5,-8.5865) (7.75,-8.3017) (8,-8.0259) (8.25,-7.7586) (8.5,-7.4993) (8.75,-7.2475) (9,-7.0029)};
\end{axis}
\end{tikzpicture}

\caption{SER dans l'axe de profils symétriques modifiés, selon l'indice $I$ (10~GHz). Traits pleins : optique physique ; tirets : loi $\pi r_{BA}$. Entre $I=8$ et $I=2$, l'écho de nez baisse d'environ 10~dB à épaisseur fixée.}\label{fig:v2-nez}
\end{figure}

À épaisseur fixée, le rayon varie comme $I^2$ : passer de $I=6$ (NACA standard) à $I=2$ divise le rayon par 9. La loi du point spéculaire prédit $-9{,}5$~dB ; l'optique physique donne $-7{,}5$~dB à toutes les épaisseurs (par exemple de $-13{,}1$ à $-20{,}6$~dBm à $t=12$~\%). Pour la première fois, la géométrie offre un levier sur la SER \emph{sans toucher à l'épaisseur}.

\begin{figure}[htbp]
\centering
% figure générée par tools/make_figures_v2.py : v2_carte
\begin{tikzpicture}
\begin{axis}[width=0.5\linewidth,height=6cm,xlabel={indice $I$},ylabel={épaisseur $t$ (\%)},
  xtick={0,1,2,3},xticklabels={2,4,6,8},ytick={6,8,...,20},enlargelimits=false,axis on top,
  colormap={bl}{color=(white) color=(cblue!45) color=(cblue!85!black)},colorbar,
  colorbar style={ylabel={SER S1 (dBm)},width=2.5mm,tick label style={font=\scriptsize}},tick label style={font=\scriptsize}]
\addplot[matrix plot*,mesh/cols=4,point meta=explicit] coordinates {(0,6) [-26.108] (1,6) [-21.571] (2,6) [-18.655] (3,6) [-16.578] (0,8) [-23.454] (1,8) [-19.088] (2,8) [-16.222] (3,8) [-14.169] (0,10) [-21.462] (1,10) [-17.102] (2,10) [-14.306] (3,10) [-12.353] (0,12) [-19.794] (1,12) [-15.477] (2,12) [-12.764] (3,12) [-10.865] (0,14) [-18.364] (1,14) [-14.096] (2,14) [-11.470] (3,14) [-9.641] (0,16) [-17.106] (1,16) [-12.896] (2,16) [-10.362] (3,16) [-8.624] (0,18) [-16.010] (1,18) [-11.856] (2,18) [-9.396] (3,18) [-7.712] (0,20) [-15.023] (1,20) [-10.930] (2,20) [-8.512] (3,20) [-6.897]};
\end{axis}
\end{tikzpicture}

\caption{SER moyenne du secteur S1 sur la grille V2 (médiane sur la cambrure et sa position, à l'incidence de croisière). L'épaisseur et le rayon de bord d'attaque jouent désormais des rôles comparables (corrélations 0,71 et 0,67).}\label{fig:v2-carte}
\end{figure}

\section{Attention : les nez aigus sortent du régime optique}
Le prix de ce levier est un problème de validité. Pour $I=2$, à 10~GHz :
\begin{center}
\begin{tabular}{lcc}
\toprule
 & $t=8$~\% & $t=12$~\% \\
\midrule
$r_{BA}$ & 0,8~mm & 1,8~mm \\
$k\,r_{BA}$ & 0,16 & 0,37 \\
écart Mie / PO pour un cylindre (TM) & $+4{,}6$~dB & $+2{,}5$~dB \\
\bottomrule
\end{tabular}
\end{center}
Le nez est alors bien plus petit que la longueur d'onde : on n'est plus en régime optique, mais dans la zone de résonance, voire de Rayleigh. La solution exacte du cylindre montre qu'en polarisation TM, l'écho réel est \emph{plus fort} que $\pi a$, de 2,5 à 4,6~dB. Le gain de SER d'un nez aigu prédit par l'optique physique est donc probablement surestimé, peut-être de moitié en TM ; en TE, l'écart est de signe opposé.

\section{La réponse exacte : méthode des moments}
\begin{evolution}{V2}{section ajoutée après le premier calcul de la V2}Pour lever le doute, on a écrit une méthode des moments 2D en polarisation TM (\texttt{Code/ser/mom2d.py}), c'est-à-dire la résolution exacte de l'équation intégrale du champ électrique sur le contour. Elle reproduit la série de Mie du cylindre à 0,02~dB près, de $ka=0{,}5$ à $ka=31$, et ses résultats sur les profils sont convergés en maillage ($\lambda/15$, raffiné à $r_{BA}/4$ au nez) à 0,02~dB.\end{evolution}

\begin{figure}[htbp]
\centering
% figure générée par tools/make_figures_v2.py : v2_mom
\begin{tikzpicture}
\begin{axis}[width=0.72\linewidth,height=5.8cm,xlabel={indice de rayon de bord d'attaque $I$},ylabel={SER moyenne S1 (dBm)},
  xmin=1.5,xmax=8.5,legend style={draw=none,font=\scriptsize},legend pos=south east]
\addplot[corange,very thick,mark=*,mark size=1.4pt] coordinates {(2,-19.455) (3,-18.361) (4,-17.323) (6,-15.44) (8,-13.823)};\addlegendentry{MoM, $t=8$~\%}
\addplot[corange,dashed,mark=o,mark size=1.4pt,forget plot] coordinates {(2,-23.76) (3,-21.275) (4,-19.321) (6,-16.401) (8,-14.27)};
\addplot[cblue,very thick,mark=*,mark size=1.4pt] coordinates {(2,-17.679) (3,-16.191) (4,-14.846) (6,-12.585) (8,-10.775)};\addlegendentry{MoM, $t=12$~\%}
\addplot[cblue,dashed,mark=o,mark size=1.4pt,forget plot] coordinates {(2,-20.159) (3,-17.68) (4,-15.753) (6,-12.905) (8,-10.876)};
\addplot[cgreen,very thick,mark=*,mark size=1.4pt] coordinates {(2,-16.001) (3,-14.255) (4,-12.739) (6,-10.322) (8,-8.5093)};\addlegendentry{MoM, $t=16$~\%}
\addplot[cgreen,dashed,mark=o,mark size=1.4pt,forget plot] coordinates {(2,-17.545) (3,-15.083) (4,-13.187) (6,-10.441) (8,-8.5384)};
\end{axis}
\end{tikzpicture}

\caption{SER moyenne du secteur S1 de profils symétriques : méthode des moments (traits pleins) et optique physique (tirets). Les deux méthodes coïncident pour les nez ronds ; pour les nez aigus, l'optique physique sous-estime l'écho de 2,5 à 4,3~dB.}\label{fig:v2-mom}
\end{figure}

\paragraph{Ce que change la méthode exacte.}
\begin{itemize}
\item Le \textbf{gain d'un nez aigu existe}, mais il est plus faible : de $I=6$ à $I=2$, $-5{,}1$~dB à $t=12$~\% (au lieu de $-7{,}3$~dB) et $-4{,}0$~dB à $t=8$~\% (au lieu de $-7{,}4$~dB).
\item Sur les 1\,024 profils, l'optique physique sous-estime la SER de 1,05~dB en moyenne (écart-type 1,3~dB), mais le \textbf{classement} est presque identique : corrélation de rang 0,999.
\item Les \textbf{fronts de Pareto} sont robustes : 3 profils communs sur 4 dans S1, 9 sur 10 dans S2. Seul l'écart de SER le long du front diminue : de 1,8 à 0,8~dB dans S1, de 5,3 à 3,4~dB dans S2.
\end{itemize}

\begin{projet}
La méthode des moments devient la SER de référence de la V2 (\code{with\_objective(..., source="MoM")}) ; l'optique physique reste l'outil rapide. Le compromis traînée/SER de la V2 survit à la méthode exacte, avec un bénéfice radar plus modeste. Il reste à vérifier la polarisation TE et le passage en 3D, ce que le solveur du groupe SER, ou \texttt{bempp-cl} sur le STL, pourra faire.
\end{projet}

\exercices
\begin{exercice}
À quelle fréquence le nez d'un profil $t=12$~\%, $I=2$ (corde 1~m) atteindrait-il $k\,r_{BA}=3$, où l'optique physique est à 0,3~dB près ? Quelle est la longueur d'onde correspondante ?
\end{exercice}

\begin{pistes}
\begin{itemize}
\item \iterbadge{V2}~\textbf{Fait} : méthode des moments 2D en TM. Reste la polarisation TE (équation intégrale du champ magnétique), plus délicate.
\item Calculer la SER dans les deux polarisations et retenir la pire, ce qui est l'usage pour une menace de polarisation inconnue.
\end{itemize}
\end{pistes}


\appendix
\chapter{Corrigés des exercices}\label{annexe:corriges}
\iterbanner{V1}

\section*{Chapitre~\ref{ch:intro}}
\textbf{Ex.~1.1.} $\sigma_{2D}=10^{-1{,}3}=0{,}050$~m. $\sigma_{3D}=kb^2\sigma_{2D}/\pi=209{,}6\times4\times0{,}050/\pi\approx13$~m\textsuperscript{2}, soit environ 11~dBsm. C'est une valeur énorme parce que la direction perpendiculaire à l'envergure est un \emph{spéculaire de ligne} : tout le bord d'attaque renvoie en phase. Un avion furtif évite précisément cette configuration (flèche du bord d'attaque, qui renvoie l'écho hors de l'axe de la menace).

\textbf{Ex.~1.2.} $k=209{,}6$~m\textsuperscript{-1}. $r_{BA}=4{,}0$, $15{,}9$ et $44{,}1$~mm, soit $kr_{BA}=0{,}83$, $3{,}3$ et $9{,}2$. Pour $kr\ge5$ avec $r=4$~mm : $k\ge1\,250$~m\textsuperscript{-1}, soit $f\ge60$~GHz.

\section*{Chapitre~\ref{ch:mecanismes}}
\textbf{Ex.~2.1.} Facette : $kw^2=209{,}6\times0{,}01=2{,}1$~m, soit $+3{,}2$~dBm. Bord d'attaque : $\pi\times0{,}016=0{,}050$~m, soit $-13$~dBm. Écart : environ 16~dB. Inclinée de $10^\circ$ : $kw\sin10^\circ=3{,}6$, $\operatorname{sinc}^2(3{,}6)\approx0{,}015$, soit $-18$~dB par rapport au pic : la facette tombe vers $-15$~dBm. Une facette est très forte dans sa direction et très faible ailleurs : c'est le principe du F-117.

\section*{Chapitre~\ref{ch:po}}
\textbf{Ex.~3.1.} Les deux facettes ont des normales opposées ; une seule est éclairée ($\vect n\cdot\shat>0$). Pour elle, $\vect n\cdot\shat=\cos\theta$, $\vect t\cdot\shat=\sin\theta$, $L=w$ : $\sigma_{2D}=k\,|\cos\theta\,w\,\operatorname{sinc}(kw\sin\theta)|^2$, la formule de la bande.

\textbf{Ex.~3.2.} $\sigma_{2D}=\pi\rho=\pi b^2/a$. Avec $a=0{,}5$~m et $b=0{,}06$~m : $\pi\times0{,}0072=0{,}023$~m, soit $-16{,}4$~dBm, environ 3~dB de moins qu'un NACA 0012 de même épaisseur : l'ellipse a un nez plus pointu ($\rho=7{,}2$~mm contre 15,9~mm).

\section*{Chapitre~\ref{ch:code}}
\textbf{Ex.~4.1.} Environ $400\times81\approx3\cdot10^4$ évaluations de sinus, cosinus et exponentielles, soit quelques dizaines de microsecondes de calcul effectif ; le reste est la construction du contour. Une méthode des moments à 400 inconnues demande de remplir une matrice de $1{,}6\cdot10^5$ termes (fonctions de Hankel) puis une factorisation en $N^3/3\approx2\cdot10^7$ opérations, mais pour \emph{toutes} les directions à la fois : quelques dizaines de millisecondes. C'est encore très raisonnable en 2D.

\section*{Chapitre~\ref{ch:naca}}
\textbf{Ex.~5.1.} $20\log_{10}(8/12)=-3{,}5$~dB. Côté aérodynamique, à $\alpha=2^\circ$, la traînée diminue aussi quand l'épaisseur baisse (quelques counts entre 12 et 8~\% pour un profil symétrique, voir la base). Les deux objectifs vont dans le même sens : c'est l'origine du front dégénéré.

\textbf{Ex.~5.2.} Paramétrons $x=y^2/(2r)$. Alors $x'=y/r$, $x''=1/r$, et la courbure vaut $\kappa=|x''|/(1+x'^2)^{3/2}=\frac{1/r}{(1+y^2/r^2)^{3/2}}$. La normale fait l'angle $\beta$ avec l'axe des $x$ quand la pente $\dd x/\dd y=y/r=\tan\beta$. Donc $1+y^2/r^2=1/\cos^2\beta$ et $\rho=1/\kappa=r/\cos^3\beta$.

\section*{Chapitre~\ref{ch:integration}}
\textbf{Ex.~6.1.} Dans le repère du profil, $\shat=-(\cos(\alpha-\varepsilon),\sin(\alpha-\varepsilon))$. Pour $\alpha=-4^\circ$ et $\varepsilon=0$, $\shat$ pointe vers l'avant et vers le \emph{haut} du profil : le radar voit davantage l'extrados. À $+2^\circ$, il le voit par en dessous, côté intrados. L'extrados, plus bombé, présente au radar une courbure différente : l'écho dans le secteur se déplace, et la valeur moyenne change de quelques dixièmes de décibel à quelques décibels selon le secteur.

\section*{Chapitre~\ref{ch:limites}}
\textbf{Ex.~7.1.} En TM, le champ électrique est parallèle à l'envergure : l'équation intégrale est scalaire et son noyau est simplement $H_0^{(2)}$, avec une singularité logarithmique facile à intégrer. En TE, il faut la dérivée normale du noyau, plus singulière. Il faut environ 10 inconnues par longueur d'onde pour représenter correctement la variation du courant (deux fois Nyquist, avec une marge). Le contour d'un profil mesure environ $2c=2$~m, soit 67 longueurs d'onde à 10~GHz : environ 700 inconnues, plus un raffinement au bord d'attaque.

\section*{Chapitre~\ref{ch:v2} \iterbadge{V2}}
\textbf{Ex.~8.1.} $r_{BA}=1{,}1019\times(0{,}12\times2/6)^2=1{,}76$~mm. $k=3/r=1\,700$~m\textsuperscript{-1}, soit $f=kc_0/(2\pi)\approx81$~GHz et $\lambda\approx3{,}7$~mm. Aucun radar de surveillance usuel n'opère à cette fréquence : pour ces nez, il faut un solveur exact.


\end{document}
